\phantomsection
\part{線型代数の基礎}

\begin{abstract}
  この章では，線型代数の基礎となる用語や概念を整理する．
  まず行列を定義し，その和・スカラー倍・積といった計算規則と逆行列について確認する．
  次に，行列を用いた連立一次方程式の取り扱いを見る．
  そのうえで線型写像を定義し，$K^n$から$K^m$への線型写像が$m \times n$行列とちょうど一対一に対応することを示す．
  また，線型写像に付随する基本的な集合として，像と核を導入する．
  最後に線型結合を定義し，「線型独立」「線型従属」の概念および行列の階数を導入する．
\end{abstract}

\phantomsection
\section{行列}

\phantomsection
\subsection{行列の定義}

以下，本書では$K$は$\mathbb{R}$または$\mathbb{C}$を表すと約束する．
また，本書では自然数に$0$を含めない．すなわち，自然数とは正の整数$1, 2, 3, \dots$のことであり，自然数全体の集合を$\mathbb{N} = \{1, 2, 3, \dots\}$と表す．

まず，本書を通じて主役となる\textbf{行列}\index{ぎょうれつ@行列}を定義しよう．

\begin{definition}{行列}{行列}
  $mn$個の数$a_{ij} \in K$~$(1 \le i \le m,~1 \le j \le n)$を，縦に$m$個，横に$n$個並べた表
  \[
    A =
    \begin{pmatrix}
      a_{11} & a_{12} & \cdots & a_{1n} \\
      a_{21} & a_{22} & \cdots & a_{2n} \\
      \vdots & \vdots & \ddots & \vdots \\
      a_{m1} & a_{m2} & \cdots & a_{mn}
    \end{pmatrix}
  \]
  を\textbf{$m \times n$行列}（$m \times n$型の行列）という．
  横の並びを\textbf{行}\index{ぎょう@行}，縦の並びを\textbf{列}\index{れつ@列}とよび，第$i$行と第$j$列の交わる位置にある数$a_{ij}$を$A$の\textbf{$(i, j)$成分}\index{せいぶん@成分}という．
  また，この$A$を$A = (a_{ij})$と略記することがある．
\end{definition}

高校までの学習で，列ベクトル
\[
  \begin{pmatrix} x_1 \\ x_2 \\ \vdots \\ x_n \end{pmatrix}
\]
を学ぶが，これは$n \times 1$行列である．同様に行ベクトル
\[
  \begin{pmatrix} x_1 & x_2 & \cdots & x_n \end{pmatrix}
\]
は$1 \times n$行列である．なお，以下では行ベクトルを
\[
  (x_1, x_2, \dots, x_n)
\]
のようにカンマ区切りで表す表記も用いることにする．

\begin{example}{行列}{行列}
  \[
    A =
    \begin{pmatrix}
      -1 & 5 & 3 & 4 \\
      9 & -2 & 4 & 1 \\
      6 & -10 & -3 & 5
    \end{pmatrix}
  \]
  としたとき，$A$は$3 \times 4$行列である．
  \[
    \begin{pmatrix} -1 & 5 & 3 & 4 \end{pmatrix}
  \]
  を第$1$行といい，同様に第$2$行，第$3$行が定まる．
  また，
  \[
    \begin{pmatrix} -1 \\ 9 \\ 6 \end{pmatrix}
  \]
  を第$1$列といい，同様に第$2$列から第$4$列までが定まる．
  たとえば，$A$の$(2, 3)$成分は$4$である．
\end{example}

新しい対象を定義したら，次に「$2$つの対象がいつ等しいとみなされるのか」を約束しておく必要がある．
行列については，次のように定める．
 
\begin{definition}{行列の相等}{行列の相等}
  $2$つの行列$A = (a_{ij})$，$B = (b_{ij})$が\textbf{等しい}\index{ぎょうれつのそうとう@行列の相等}とは，$A$と$B$の型が等しく，かつ対応するすべての成分が等しいこと，すなわち$A$，$B$がともに$m \times n$行列であって
  \[
    a_{ij} = b_{ij} \quad (1 \le i \le m,~1 \le j \le n)
  \]
  が成り立つことをいう．このとき$A = B$とかく．
\end{definition}
 
型が等しいことも条件に含まれている点に注意してほしい．
たとえば，$2 \times 2$の零行列と$2 \times 3$の零行列は，どちらも$O$という同じ記号で表されるが，等しい行列ではない．
 
また，この定義によって，$1$つの行列の等式$A = B$は$mn$個の数の等式$a_{ij} = b_{ij}$を束ねたものとなる．
のちに逆行列を求める場面などで，行列の等式から成分を比較して連立方程式を取り出す議論が現れるが，それを支えているのがこの定義である．
 
\begin{mycolumn}
  「等しい」という関係をわざわざ定義することを，不思議に思う方もいるかもしれない．
  しかし，これは数学の議論を始めるうえで欠かせない約束である．
 
  実数の世界を思い出してみよう．円周率$\pi$は，小数表示$3.1415\ldots$のほかに，無限級数
  \[
    4 \left( 1 - \frac{1}{3} + \frac{1}{5} - \frac{1}{7} + \cdots \right)
  \]
  としても，積分
  \[
    \int_0^1 \frac{4}{1 + x^2} \, dx
  \]
  としてもかける．見た目のまったく異なるこれらの式を「$=$」で結べるのは，実数における「$=$」が「表記のかたちが同じ」ではなく「同じ数を表している」という意味に定められているからである．
 
  行列でも事情は同じで，$1$つの行列がさまざまな計算の結果として現れうる．
  そこで，どのような経緯で得られた行列であっても，「型が等しく，すべての成分が等しければ同じ行列とみなす」と約束しておくのである．
  逆にいえば，この約束があるからこそ，これから定義する和や積の計算結果を$1$つの行列として自信をもって「$=$」で結ぶことができる．
\end{mycolumn}

\begin{definition}{正方行列}{正方行列}
  行の数と列の数が等しい行列を\textbf{正方行列}\index{せいほうぎょうれつ@正方行列}という．
  すなわち，$n$を自然数として，$n \times n$行列を\textbf{$n$次正方行列}という．
\end{definition}

\begin{example}{正方行列}{正方行列}
  \[
    \begin{pmatrix} 2 & 5 \\ -8 & 7 \end{pmatrix}, \quad
    \begin{pmatrix} -1 & 4 & 0 \\ 2 & 1 & 9 \\ -8 & 5 & 1 \end{pmatrix}
  \]
  はそれぞれ$2$次，$3$次の正方行列である．
\end{example}

行列全体の集合にも記号を用意しておこう．

\begin{definition}{行列全体の集合}{行列全体の集合}
  $K$の元\footnote{集合に属する個々の対象を，その集合の\textbf{元}（\textbf{要素}）とよぶ．本書では「元」の語に統一する．}を成分とする$m \times n$行列全体の集合を
  \[
    M_{m,n}(K)
  \]
  と表す．とくに，$n$次正方行列全体の集合を
  \[
    M_n(K) = M_{n,n}(K)
  \]
  と表す．
\end{definition}

この記法は，のちに「行列を変数とする写像」を考える場面（たとえば行列式）や，「$m \times n$行列全体が一つの線型空間をなす」ことを述べる場面（第3章）で用いられる．

のちの議論で頻繁に現れる「特別な行列」にも名前をつけておこう．

\begin{definition}{零行列と単位行列}{零行列と単位行列}
  すべての成分が$0$である行列を\textbf{零行列}\index{れいぎょうれつ@零行列}とよび，$O$と表す．

  また，$(1,1)$成分から$(n,n)$成分までの対角線上の成分がすべて$1$で，他の成分がすべて$0$である$n$次正方行列
  \[
    E_n =
    \begin{pmatrix}
      1 & 0 & \cdots & 0 \\
      0 & 1 & \cdots & 0 \\
      \vdots & \vdots & \ddots & \vdots \\
      0 & 0 & \cdots & 1
    \end{pmatrix}
  \]
  を\textbf{$n$次単位行列}\index{たんいぎょうれつ@単位行列}とよぶ\footnote{単位行列を表す文字は$E$のほかに$I$も使われる．英語表記の``Identity Matrix''の頭文字の$I$をとるか，もしくはドイツ語表記の``Einheitsmatrix''の頭文字をとるかの選択である．}．
  サイズが文脈から明らかなときは，単に$E$とかく．
\end{definition}

\begin{definition}{転置行列}{転置行列}
  $m \times n$行列$A = (a_{ij})$に対して，$A$の行と列を入れ替えて得られる$n \times m$行列，すなわち$(i, j)$成分が$a_{ji}$である行列を$A$の\textbf{転置行列}\index{てんちぎょうれつ@転置行列}とよび，$A^\mathsf{T}$と表す\footnote{行列$A$の転置行列の表記として，$A^\mathsf{T}$のほかに${}^t A$もよく使われる．}．
\end{definition}

\begin{example}{転置行列}{転置行列}
  \[
    A = \begin{pmatrix} 1 & 2 & 3 \\ 4 & 5 & 6 \end{pmatrix}
    \quad \text{のとき，} \quad
    A^\mathsf{T} = \begin{pmatrix} 1 & 4 \\ 2 & 5 \\ 3 & 6 \end{pmatrix}
  \]
  である．$A$の第$1$行が$A^\mathsf{T}$の第$1$列になっていることに注意してほしい．
\end{example}

いきなり行列という概念が出てきて，戸惑っている方も多いと思う．
これから行列の計算規則や意味などを確認していこう．

\phantomsection
\subsection{行列の和とスカラー倍}

\begin{definition}{行列の和}{行列の和}
  $m \times n$行列$A = (a_{ij})$，$B = (b_{ij})$に対して，その\textbf{和}\index{ぎょうれつのわ@行列の和}$A + B$を
  \[
    A + B =
    \begin{pmatrix}
      a_{11} + b_{11} & a_{12} + b_{12} & \cdots & a_{1n} + b_{1n} \\
      a_{21} + b_{21} & a_{22} + b_{22} & \cdots & a_{2n} + b_{2n} \\
      \vdots & \vdots & \ddots & \vdots \\
      a_{m1} + b_{m1} & a_{m2} + b_{m2} & \cdots & a_{mn} + b_{mn}
    \end{pmatrix}
  \]
  と定義する．すなわち，成分ごとに和をとる．
\end{definition}

型の異なる行列に対しては，和は定義されない．

\begin{definition}{行列のスカラー倍}{行列のスカラー倍}
  $\alpha \in K$と$m \times n$行列$A = (a_{ij})$に対して，$A$の\textbf{スカラー倍}\index{ぎょうれつのすからーばい@行列のスカラー倍}$\alpha A$を
  \[
    \alpha A =
    \begin{pmatrix}
      \alpha a_{11} & \alpha a_{12} & \cdots & \alpha a_{1n} \\
      \alpha a_{21} & \alpha a_{22} & \cdots & \alpha a_{2n} \\
      \vdots & \vdots & \ddots & \vdots \\
      \alpha a_{m1} & \alpha a_{m2} & \cdots & \alpha a_{mn}
    \end{pmatrix}
  \]
  と定義する．すなわち，すべての成分を$\alpha$倍する．
\end{definition}

\begin{example}{行列のスカラー倍}{行列のスカラー倍}
  $A = \begin{pmatrix} 2 & 3 \\ 9 & 1 \\ -8 & 3 \end{pmatrix}$に$5$をかけると，
  \[
    5A =
    \begin{pmatrix}
      5 \cdot 2 & 5 \cdot 3 \\
      5 \cdot 9 & 5 \cdot 1 \\
      5 \cdot (-8) & 5 \cdot 3
    \end{pmatrix}
    =
    \begin{pmatrix}
      10 & 15 \\
      45 & 5 \\
      -40 & 15
    \end{pmatrix}
  \]
  となる．
\end{example}

\begin{mycolumn}
  行列のスカラー倍を，たとえば「第$1$行だけを$\alpha$倍する」操作として定めることも，定義としては可能である．
  この操作を仮に$\alpha \ast A$とかくと，$1 \ast A = A$や$\alpha \ast (A + B) = \alpha \ast A + \alpha \ast B$は成り立つが，
  \[
    (\alpha + \beta) \ast A = \alpha \ast A + \beta \ast A
  \]
  は一般に成り立たない．実際，左辺の第$2$行以降は$A$のままであるのに対し，右辺の第$2$行以降は$A$の対応する行の$2$倍になる．
  \deref{行列のスカラー倍}が「すべての成分を$\alpha$倍する」という形をしているのは，こうした計算規則が成り立つように選ばれた結果である．
  定義は天下りに与えられるものではなく，望ましい性質から逆算して設計されるものである．
\end{mycolumn}

\begin{definition}{行列の減法}{行列の減法}
  $A$，$B$を型の等しい行列とする．
  \deref{行列のスカラー倍}に従って行列$(-1)B$を考え，
  \[
    A + (-1)B
  \]
  を$A - B$とかく．
  以下，$(-1)B$を$-B$と略記する．
\end{definition}

\deref{行列の減法}は，減法を和とスカラー倍から組み立てる定義である．
実際に成分を計算すれば，$A - B$の$(i, j)$成分は$a_{ij} - b_{ij}$であり，成分ごとに差をとることと一致する．

\begin{example}{行列の和と差の計算}{行列の和と差の計算}
  \[
    A = \begin{pmatrix} 2 & 0 & -1 \\ 1 & 1 & 4 \end{pmatrix}, \quad
    B = \begin{pmatrix} 1 & 2 & 3 \\ 0 & 1 & -2 \end{pmatrix}, \quad
    C = \begin{pmatrix} 0 & 1 \\ 3 & 5 \\ 2 & 0 \end{pmatrix}
  \]
  としたとき，
  \[
    A + B =
    \begin{pmatrix}
      2 + 1 & 0 + 2 & -1 + 3 \\
      1 + 0 & 1 + 1 & 4 + (-2)
    \end{pmatrix}
    =
    \begin{pmatrix}
      3 & 2 & 2 \\
      1 & 2 & 2
    \end{pmatrix}, \quad
    A - B =
    \begin{pmatrix}
      1 & -2 & -4 \\
      1 & 0 & 6
    \end{pmatrix}
  \]
  である．一方，$A$と$C$は型が異なるので，$A + C$や$A - C$は定義されない．
\end{example}

行列の和について成り立つ計算規則をまとめておこう．

\begin{prop}{行列の和の計算規則}{行列の和の計算規則}
  $A$，$B$，$C$を型の等しい行列とし，$O$を同じ型の零行列とする．このとき，次が成り立つ：
  \begin{align*}
    (A + B) + C &= A + (B + C) && \text{（和の結合律）} \\
    A + B &= B + A && \text{（和の交換律）} \\
    A + O &= A && \text{（加法単位元の存在）} \\
    A + (-A) &= O && \text{（加法逆元の存在）}
  \end{align*}
\end{prop}

\begin{leftbar}
\begin{proof}
  行列の和とスカラー倍は成分ごとの計算として定義されているから，いずれの等式も，両辺の$(i, j)$成分を比較すれば，数$a_{ij}, b_{ij}, c_{ij} \in K$についての対応する等式に帰着される．
  たとえば第$1$式については，両辺の$(i, j)$成分はともに$a_{ij} + b_{ij} + c_{ij}$であり，これは数の和の結合律にほかならない．
  他の式も同様である．
\end{proof}
\end{leftbar}

つまり，行列の加法については，両辺が定義されるかぎり，結合律および交換律が成り立つ．
次節では行列の積を定義し，積についてもこれらの性質が成り立つかを考えることにする．

\phantomsection
\subsection{行列の積}

\begin{definition}{行列の積}{行列の積}
  $l \times m$行列$A = (a_{ij})$と$m \times n$行列$B = (b_{jk})$に対して，$A$と$B$の\textbf{積}\index{ぎょうれつのせき@行列の積}$AB$を，$(i, k)$成分が
  \[
    c_{ik} = \sum_{j=1}^{m} a_{ij} b_{jk}
    \quad (i = 1, 2, \dots, l,~ k = 1, 2, \dots, n)
  \]
  で与えられる$l \times n$行列$C = (c_{ik})$として定義する．
\end{definition}


\begin{wrapfigure}[7]{r}{0.5\textwidth}
  \centering
  \vspace{-\baselineskip}
  \begin{tikzpicture}[
      mat/.style={matrix of math nodes,
                  left delimiter=(, right delimiter=),
                  nodes={inner sep=1.5pt, font=\scriptsize},
                  column sep=5pt, row sep=4pt},
      waku/.style={draw, magenta, thick, rounded corners=2pt, inner sep=2.5pt}]
    \matrix (A) [mat] {
      a_{11} & \cdots & a_{1m} \\
      \vdots &        & \vdots \\
      |[magenta]| a_{i1} & |[magenta]| \cdots & |[magenta]| a_{im} \\
      a_{l1} & \cdots & a_{lm} \\
    };
    \matrix (B) [mat, right=14pt of A] {
      b_{11} & \cdots & |[magenta]| b_{1k} & \cdots & b_{1n} \\
      \vdots &        & |[magenta]| \vdots &        & \vdots \\
      b_{m1} & \cdots & |[magenta]| b_{mk} & \cdots & b_{mn} \\
    };
    \node[waku, fit=(A-3-1)(A-3-3)] (row) {};
    \node[waku, fit=(B-1-3)(B-3-3)] (col) {};
    \node[magenta, font=\tiny, left=8pt of row] {第$i$行};
    \node[magenta, font=\tiny, above=4pt of col] {第$k$列};
  \end{tikzpicture}
    \caption{行列の積の計算}
  \label{fig:matrix-product}
  \vspace{-.5\baselineskip}
\end{wrapfigure}

積$AB$が定義されるのは，左の行列$A$の列の数と，右の行列$B$の行の数が一致するときに限ることに注意したい．
そして，$AB$の$(i, k)$成分は，$A$の第$i$行と$B$の第$k$列の成分どうしを順にかけて足し合わせたもの，すなわち
\[
  a_{i1}b_{1k} + a_{i2}b_{2k} + \cdots + a_{im}b_{mk}
\]
である．

\deref{行列の積}は複雑である．
しかし，のちに線型写像について考察すると，この定義が自然であることが明らかとなる\footnote{行列の積が写像の合成に対応することについては，\prref{行列の積と写像の合成}で確認する．}．
ここでは具体例で確認してみよう．

\begin{example}{行列の積}{行列の積}
  \[
    A = \begin{pmatrix} 1 & 2 & 0 \\ 3 & 2 & 0 \end{pmatrix}, \quad
    B = \begin{pmatrix} 1 & 0 \\ 0 & 2 \\ 1 & 1 \end{pmatrix}
  \]
  のとき，$A$は$2 \times 3$行列，$B$は$3 \times 2$行列であるため，積$AB$は定義され，
  \[
    AB =
    \begin{pmatrix}
      1 \cdot 1 + 2 \cdot 0 + 0 \cdot 1 & 1 \cdot 0 + 2 \cdot 2 + 0 \cdot 1 \\
      3 \cdot 1 + 2 \cdot 0 + 0 \cdot 1 & 3 \cdot 0 + 2 \cdot 2 + 0 \cdot 1
    \end{pmatrix}
    =
    \begin{pmatrix}
      1 & 4 \\
      3 & 4
    \end{pmatrix}
  \]
  である．
\end{example}

一般に，積$AB$が定義されても$BA$が定義されるとは限らない．
ただし，$A$が$m\times n$行列，$B$が$n\times m$行列であれば，
$AB$と$BA$はともに定義される．
とくに，同じ次数の正方行列どうしでは，常に$AB$と$BA$の両方が定義される．

任意の$a, b \in \mathbb{R}$に対して，以下の$2$つの事実が成り立つことは既知であろう：
\begin{gather*}
  ab = ba, \\
  ab = 0 \text{ ならば } a = 0 \text{ または } b = 0.
\end{gather*}
つまり，$\mathbb{R}$では積の交換律が成り立ち，$a \ne 0$かつ$b \ne 0$なのに$ab = 0$となるような$a, b$は存在しない．
では，行列の場合はどうであろうか．

\begin{example}{行列の積の非可換性と零因子}{行列の積の非可換性と零因子}
  行列$A$，$B$を
  \[
    A = \begin{pmatrix} 1 & 0 \\ 1 & 0 \end{pmatrix}, \quad
    B = \begin{pmatrix} 0 & 0 \\ 0 & 1 \end{pmatrix}
  \]
  で定めると，
  \[
    AB = \begin{pmatrix} 0 & 0 \\ 0 & 0 \end{pmatrix}, \quad
    BA = \begin{pmatrix} 0 & 0 \\ 1 & 0 \end{pmatrix}
  \]
  となる．
\end{example}

この例から，行列においては，一般に積の交換律$AB = BA$は成立しないことがわかる．
さらに，$A \ne O$かつ$B \ne O$であるのに$AB = O$となることも起こりうる\footnote{このように，零行列でないのに，零行列でない行列をかけると積が零行列になるような行列を\textbf{零因子}\index{れいいんし@零因子}という．}．
このように，行列の積は数の積とは異なる振る舞いをすることがあるので，計算の際には注意が必要である．

一方で，行列の積は数の積と似た性質をもつ場面もある．

\begin{prop}{行列の乗法結合律}{行列の乗法結合律}
  $A$が$l \times m$行列，$B$が$m \times n$行列，$C$が$n \times p$行列であれば，
  \[
    (AB)C = A(BC)
  \]
  が成り立つ．
\end{prop}

\begin{leftbar}
\begin{proof}
  $A = (a_{ij})$，$B = (b_{jk})$，$C = (c_{kq})$とかく．
  $AB$の$(i, k)$成分は
  \[
    \sum_{j=1}^{m} a_{ij} b_{jk}
    \quad (i = 1, 2, \dots, l,~ k = 1, 2, \dots, n)
  \]
  であるから，$(AB)C$の$(i, q)$成分は
  \[
    \sum_{k=1}^{n} \left( \sum_{j=1}^{m} a_{ij} b_{jk} \right) c_{kq}
    = \sum_{k=1}^{n} \sum_{j=1}^{m} a_{ij} b_{jk} c_{kq}
  \]
  である．同様に，$A(BC)$の$(i, q)$成分は
  \[
    \sum_{j=1}^{m} \sum_{k=1}^{n} a_{ij} b_{jk} c_{kq}
  \]
  である．この$2$つの式は，$\sum$記号の順序が違うだけで一致している．
  これが証明すべきことであった．
\end{proof}
\end{leftbar}

交換律は成り立たないが結合律は成り立つ，という場面はほかにもある．
たとえば，写像の合成について，
\[
  f \circ g = g \circ f
\]
は一般には成り立たないが，
\[
  h \circ (g \circ f) = (h \circ g) \circ f
\]
はつねに成り立つ．

\begin{example}{合成関数と交換律}{合成関数と交換律}
  実数全体で定義された関数$f(x) = 1 + x$，$g(x) = x^3$を考える．
  $f$と$g$，$g$と$f$はいずれも合成可能で，
  \[
    (g \circ f)(x) = (1 + x)^3, \quad (f \circ g)(x) = 1 + x^3
  \]
  となり，$g \circ f \ne f \circ g$である．
\end{example}

実は，行列の積と写像の合成のこの類似は偶然ではない．
このことは，のちに線型写像と行列の関係を調べる際に明らかになる（\prref{行列の積と写像の合成}）．

行列の積について成り立つ計算規則をまとめておこう．

\begin{prop}{行列の積の計算規則}{行列の積の計算規則}
  $A$，$B$，$C$を行列とし，$E$を単位行列とする．両辺が定義されるかぎり，次が成り立つ：
  \begin{align*}
    (AB)C &= A(BC) && \text{（積の結合律）} \\
    AE &= A, \quad EA = A && \text{（乗法単位元の存在）} \\
    A(B + C) &= AB + AC, \quad (A + B)C = AC + BC && \text{（分配律）}
  \end{align*}
\end{prop}

\begin{leftbar}
\begin{proof}
  第$1$式は\prref{行列の乗法結合律}で示した通りである．
  第$2$式と第$3$式は，\prref{行列の和の計算規則}の証明と同様に，両辺の成分を比較することで確かめられる．
\end{proof}
\end{leftbar}

\phantomsection
\subsection{逆行列}

数$a \ne 0$に対しては，$ax = xa = 1$を満たす数$x = a^{-1}$が存在した．
この「逆数」にあたる概念を，正方行列に対して考えよう．

\begin{definition}{逆行列}{逆行列}
  $n$次正方行列$A$と$n$次単位行列$E$に対して，
  \[
    AX = XA = E
  \]
  となる$n$次正方行列$X$が存在するとき，$A$は\textbf{正則}\index{せいそく@正則}であるという．
  このような$X$を$A$の\textbf{逆行列}\index{ぎゃくぎょうれつ@逆行列}とよび，$A^{-1}$とかく．
\end{definition}

\begin{prop}{逆行列の一意性}{逆行列の一意性}
  正方行列$A$に逆行列が存在するとき，それはただ一つに限る．
\end{prop}

\begin{leftbar}
\begin{proof}
  $X$，$Y$がともに$A$の逆行列であるとする．このとき，
  \[
    AY = YA = E
  \]
  が成り立つ．このもとで，
  \[
    X = XE = X(AY) = (XA)Y = EY = Y
  \]
  となるため，逆行列の一意性が示された．
\end{proof}
\end{leftbar}

\begin{mycolumn}
  逆行列は，定義から一意であることがわかった．
  ここで，「常に逆行列が一意であるならば，逆行列の定義に一意性を入れたほうが定義の情報量が多くなってよいはず」と考える方もおられるだろう．

  しかし，数学においては「定義は最小限の情報におさえて，その他のことは命題として示す」という慣習がある．
  定義を最小限にしたほうが，「逆行列が一意であることは，定義から証明できる事実である」ということが明確になるのである．
\end{mycolumn}

\begin{example}{2次正方行列の逆行列}{2次正方行列の逆行列}
  $A = \begin{pmatrix} a & b \\ c & d \end{pmatrix}$（成分は$K$の元）に対して，$AX = E$となる行列$X$が存在する条件を求めよう．
  $X = \begin{pmatrix} x_{11} & x_{12} \\ x_{21} & x_{22} \end{pmatrix}$とすると，
  \[
    AX =
    \begin{pmatrix}
      ax_{11} + bx_{21} & ax_{12} + bx_{22} \\
      cx_{11} + dx_{21} & cx_{12} + dx_{22}
    \end{pmatrix}
  \]
  であり，これが$E$と等しくなればよいので，
  \[
    \begin{cases}
      ax_{11} + bx_{21} = 1 \\
      cx_{11} + dx_{21} = 0
    \end{cases}, \quad
    \begin{cases}
      ax_{12} + bx_{22} = 0 \\
      cx_{12} + dx_{22} = 1
    \end{cases}
  \]
  を得る．$ad - bc \ne 0$のとき，この$2$組の連立方程式は加減法によりただ一組の解をもち，
  \[
    x_{11} = \frac{d}{ad - bc}, \quad
    x_{12} = -\frac{b}{ad - bc}, \quad
    x_{21} = -\frac{c}{ad - bc}, \quad
    x_{22} = \frac{a}{ad - bc}
  \]
  となる．この$X$に対して$XA$を計算すると，たしかに$XA = E$となる．
  よって，$ad - bc \ne 0$のとき$A$は正則であり，その逆行列は
  \[
    A^{-1} = \frac{1}{ad - bc} \begin{pmatrix} d & -b \\ -c & a \end{pmatrix}
  \]
  である\footnote{逆に$ad - bc = 0$のときは$A$は正則でないが，このことは上の計算（$ad - bc \ne 0$の場合の解の構成）だけからは直ちには従わない．のちの章で行列式を導入したあとに，一般の$n$次正方行列が正則であるための必要十分条件として証明される．}．
\end{example}

\phantomsection
\section{連立一次方程式}

\phantomsection
\subsection{加減法による解法}

次のような連立一次方程式を考える：
\begin{equation}
  \begin{cases}
    4x + 3y = 8 \\
    2x + 3y = 4
  \end{cases} \label{eq:連立一次方程式の例}
\end{equation}

初等的な解法を確認しておこう．
第$1$式から第$2$式を辺々引くと$2x = 4$を得るので，$x = 2$である．
これを第$2$式に代入すると$4 + 3y = 4$，すなわち$3y = 0$となるので，$y = 0$である．
よって，\eqref{eq:連立一次方程式の例}の解は$x = 2$，$y = 0$である．

この解法は「加減法」としてよく知られているものである．
本節では，この解法を行列のことばで表現し直すことを考える．

\phantomsection
\subsection{係数行列と行基本変形}

一般に，$m$個の式と$n$個の未知数からなる連立一次方程式
\[
  \begin{cases}
    a_{11}x_1 + a_{12}x_2 + \cdots + a_{1n}x_n = b_1 \\
    a_{21}x_1 + a_{22}x_2 + \cdots + a_{2n}x_n = b_2 \\
    \hspace{2em} \vdots \\
    a_{m1}x_1 + a_{m2}x_2 + \cdots + a_{mn}x_n = b_m
  \end{cases}
\]
は，行列の積を用いて
\[
  A\bm{x} = \bm{b}, \quad
  A = (a_{ij}), \quad
  \bm{x} = \begin{pmatrix} x_1 \\ x_2 \\ \vdots \\ x_n \end{pmatrix}, \quad
  \bm{b} = \begin{pmatrix} b_1 \\ b_2 \\ \vdots \\ b_m \end{pmatrix}
\]
とかける．

\begin{definition}{係数行列・拡大係数行列}{係数行列}
  上の連立一次方程式に対して，係数を並べた$m \times n$行列$A = (a_{ij})$を\textbf{係数行列}\index{けいすうぎょうれつ@係数行列}という．
  さらに，$A$の右側に$\bm{b}$を付け加えた$m \times (n + 1)$行列
  \[
    \begin{pmatrix}
      a_{11} & a_{12} & \cdots & a_{1n} & b_1 \\
      a_{21} & a_{22} & \cdots & a_{2n} & b_2 \\
      \vdots & \vdots & \ddots & \vdots & \vdots \\
      a_{m1} & a_{m2} & \cdots & a_{mn} & b_m
    \end{pmatrix}
  \]
  を\textbf{拡大係数行列}\index{かくだいけいすうぎょうれつ@拡大係数行列}という．
\end{definition}

たとえば，\eqref{eq:連立一次方程式の例}は
\[
  \begin{pmatrix} 4 & 3 \\ 2 & 3 \end{pmatrix}
  \begin{pmatrix} x \\ y \end{pmatrix}
  =
  \begin{pmatrix} 8 \\ 4 \end{pmatrix}
\]
とかけ，係数行列は$\begin{pmatrix} 4 & 3 \\ 2 & 3 \end{pmatrix}$，拡大係数行列は$\begin{pmatrix} 4 & 3 & 8 \\ 2 & 3 & 4 \end{pmatrix}$である．

\begin{definition}{行基本変形}{行基本変形}
  行列に対する以下の$3$種類の操作を\textbf{行基本変形}\index{ぎょうきほんへんけい@行基本変形}という．
  \begin{enumerate}[label=(\arabic*),leftmargin=3.2em]
    \item 第$i$行を$c$倍する（$c \ne 0$）\footnote{$c = 0$倍を許すと，第$i$行の情報が失われてしまい，もとの連立方程式と同値でなくなる．そのため$c \ne 0$という条件を課している．}．
    \item 第$i$行と第$j$行を入れ替える．
    \item 第$i$行に第$j$行の定数倍を加える（$i \ne j$）．
  \end{enumerate}
\end{definition}

拡大係数行列に行基本変形を施すことは，連立方程式に対して「両辺を定数倍する」「式の順序を入れ替える」「ある式に別の式の定数倍を加える」という操作を行うことに対応する．
いずれの操作でも連立方程式の解は変わらないので，行基本変形を繰り返して拡大係数行列を簡単な形に変形すれば，解が読み取れる．

たとえば，\eqref{eq:連立一次方程式の例}の拡大係数行列は，以下のように変形できる：
\begin{align*}
  \begin{pmatrix}
    4 & 3 & 8 \\
    2 & 3 & 4
  \end{pmatrix}
  &\xrightarrow{\text{第1行に第2行の$(-1)$倍を加える}}
  \begin{pmatrix}
    2 & 0 & 4 \\
    2 & 3 & 4
  \end{pmatrix} \\
  &\xrightarrow{\text{第2行に第1行の$(-1)$倍を加える}}
  \begin{pmatrix}
    2 & 0 & 4 \\
    0 & 3 & 0
  \end{pmatrix} \\
  &\xrightarrow{\text{第1行を$1/2$倍，第2行を$1/3$倍する}}
  \begin{pmatrix}
    1 & 0 & 2 \\
    0 & 1 & 0
  \end{pmatrix}
\end{align*}
最後の行列は連立方程式
\[
  \begin{cases}
    x = 2 \\
    y = 0
  \end{cases}
\]
を表しており，これがそのまま解を与えている．
この操作は，本質的には加減法による解法と同じである．

\phantomsection
\subsection{逆行列を用いた解法}

係数行列が正則な場合には，逆行列を用いて解を表すこともできる．
\eqref{eq:連立一次方程式の例}を
\begin{equation}
  A\bm{x} = \bm{b}, \quad
  A = \begin{pmatrix} 4 & 3 \\ 2 & 3 \end{pmatrix}, \quad
  \bm{x} = \begin{pmatrix} x \\ y \end{pmatrix}, \quad
  \bm{b} = \begin{pmatrix} 8 \\ 4 \end{pmatrix} \label{eq:連立一次方程式の行列表示}
\end{equation}
とかく．$4 \cdot 3 - 3 \cdot 2 = 6 \ne 0$なので，\exref{2次正方行列の逆行列}より$A$は正則であり，
\[
  A^{-1} = \frac{1}{6} \begin{pmatrix} 3 & -3 \\ -2 & 4 \end{pmatrix}
\]
である．これを\eqref{eq:連立一次方程式の行列表示}の両辺に左からかけると，
\begin{gather*}
  A^{-1}A\bm{x} = A^{-1}\bm{b} \\
  \therefore~ \bm{x} = A^{-1}\bm{b}
\end{gather*}
であり，
\[
  A^{-1}\bm{b} = \frac{1}{6}
  \begin{pmatrix} 3 & -3 \\ -2 & 4 \end{pmatrix}
  \begin{pmatrix} 8 \\ 4 \end{pmatrix}
  =
  \begin{pmatrix} 2 \\ 0 \end{pmatrix}
\]
であるから，与えられた連立一次方程式の解は$x = 2$，$y = 0$とわかる．
これは加減法および行基本変形で得た解と一致している．

\phantomsection
\section{線型写像}

\phantomsection
\subsection{写像と線型写像}

本節では，行列と密接に関わる概念である\textbf{線型写像}\index{せんけいしゃぞう@線型写像}を定義する．
その準備として，まず一般の\textbf{写像}\index{しゃぞう@写像}を定義しよう．

\begin{definition}{写像}{写像}
  $V$，$W$を集合とする．
  対応$f \colon V \to W$が\textbf{写像}であるとは，すべての$x \in V$に対して，$f(x) = y$となる$y \in W$がただひとつ存在することである．
  このとき，$V$を$f$の\textbf{定義域}\index{ていぎいき@定義域}，$W$を$f$の\textbf{終域}\index{しゅういき@終域}という．
\end{definition}

写像のうち，「和とスカラー倍を保つ」という性質をもつものを線型写像とよぶ．

\begin{definition}{線型写像・線型変換}{線型写像}
  写像$f \colon K^n \to K^m$が，任意の$\bm{x}, \bm{y} \in K^n$と任意の$\alpha \in K$に対して
  \[
    \begin{cases}
      f(\bm{x} + \bm{y}) = f(\bm{x}) + f(\bm{y}) \\
      f(\alpha\bm{x}) = \alpha f(\bm{x})
    \end{cases}
  \]
  を満たすとき，$f$を$K^n$から$K^m$への\textbf{線型写像}という．
  とくに$m = n$のとき，$f$を$K^n$の\textbf{線型変換}\index{せんけいへんかん@線型変換}という．
\end{definition}

\deref{線型写像}に現れる$2$つの条件は，次のひとつの条件にまとめることもできる：
任意の$\bm{x}, \bm{y} \in K^n$と任意の$\alpha, \beta \in K$に対して，
\[
  f(\alpha\bm{x} + \beta\bm{y}) = \alpha f(\bm{x}) + \beta f(\bm{y})
\]
が成り立つ\footnote{実際，この条件で$\alpha = \beta = 1$とすれば第$1$の条件が，$\beta = 0$とすれば第$2$の条件が得られる．逆に$2$つの条件を順に用いれば$f(\alpha\bm{x} + \beta\bm{y}) = f(\alpha\bm{x}) + f(\beta\bm{y}) = \alpha f(\bm{x}) + \beta f(\bm{y})$となる．}．

定義だけでは掴みにくいので，最も簡単な具体例を確かめておこう．

\begin{example}{線型変換}{線型変換}
  $a \in K$を定数とし，写像$f \colon K \to K$を
  \[
    f(x) = ax
  \]
  で定義する．このとき，$f$は$K$の線型変換である．
  実際，任意の$x, y \in K$と$\alpha \in K$に対して
  \begin{align*}
    f(x + y) &= a(x + y) = ax + ay = f(x) + f(y), \\
    f(\alpha x) &= a(\alpha x) = \alpha(ax) = \alpha f(x)
  \end{align*}
  が成り立つからである．
\end{example}

この例は$n = m = 1$の場合であった．
以下，$K^n$から$K^m$への線型写像の例を，いくつか具体的に見ていこう．
とくに断らない限り，本節の幾何的な例では$K = \mathbb{R}$とし，$\mathbb{R}^2$の元
\[
  \bm{x} = \begin{pmatrix} x \\ y \end{pmatrix}
\]
を座標平面上の点$(x, y)$と同一視する．

\begin{example}{恒等変換と零写像}{恒等変換と零写像}
  写像$\mathrm{id} \colon K^n \to K^n$を$\mathrm{id}(\bm{x}) = \bm{x}$で定めると，
  \[
    \mathrm{id}(\bm{x} + \bm{y}) = \bm{x} + \bm{y} = \mathrm{id}(\bm{x}) + \mathrm{id}(\bm{y}), \quad
    \mathrm{id}(\alpha\bm{x}) = \alpha\bm{x} = \alpha\,\mathrm{id}(\bm{x})
  \]
  であるから，$\mathrm{id}$は線型変換である．これを\textbf{恒等変換}\index{こうとうへんかん@恒等変換}という．

  また，すべての$\bm{x} \in K^n$に対して$f(\bm{x}) = \bm{0}$と定めた写像$f \colon K^n \to K^m$も，
  \[
    f(\bm{x} + \bm{y}) = \bm{0} = \bm{0} + \bm{0} = f(\bm{x}) + f(\bm{y}), \quad
    f(\alpha\bm{x}) = \bm{0} = \alpha\bm{0} = \alpha f(\bm{x})
  \]
  より線型写像である．これを\textbf{零写像}\index{れいしゃぞう@零写像}という．
\end{example}

\begin{example}{相似の中心が原点である拡大・縮小}{相似変換}
  $c \in K$を定数とし，写像$f \colon K^n \to K^n$を
  \[
    f(\bm{x}) = c\bm{x}
  \]
  で定める．このとき，
  \[
    f(\bm{x} + \bm{y}) = c(\bm{x} + \bm{y}) = c\bm{x} + c\bm{y} = f(\bm{x}) + f(\bm{y}), \quad
    f(\alpha\bm{x}) = c(\alpha\bm{x}) = \alpha(c\bm{x}) = \alpha f(\bm{x})
  \]
  が成り立つので，$f$は線型変換である．
  $K = \mathbb{R}$，$n = 2$のときこれは，原点を中心とする相似比$c$の拡大・縮小にあたる．
  とくに$c = 1$のときが恒等変換，$c = 0$のときが零写像，$c = -1$のときが原点に関する対称移動である．
\end{example}

次に，座標平面上の図形的な変換の例を見よう．いずれも高校数学で図形の移動として扱われるものである．

\begin{example}{対称移動と正射影}{対称移動と正射影}
  $K = \mathbb{R}$とし，写像$\mathbb{R}^2 \to \mathbb{R}^2$を次のように定める：
  \[
    f_1 \begin{pmatrix} x \\ y \end{pmatrix} = \begin{pmatrix} x \\ -y \end{pmatrix}, \quad
    f_2 \begin{pmatrix} x \\ y \end{pmatrix} = \begin{pmatrix} y \\ x \end{pmatrix}, \quad
    f_3 \begin{pmatrix} x \\ y \end{pmatrix} = \begin{pmatrix} x \\ 0 \end{pmatrix}.
  \]
  $f_1$は$x$軸に関する対称移動，$f_2$は直線$y = x$に関する対称移動，$f_3$は$x$軸への正射影（$x$軸に下ろした垂線の足への対応）である．
  これらはいずれも線型変換である．

  たとえば$f_1$について確かめよう．
  $\bm{x} = \begin{pmatrix} x_1 \\ x_2 \end{pmatrix}$，$\bm{y} = \begin{pmatrix} y_1 \\ y_2 \end{pmatrix}$，$\alpha \in \mathbb{R}$に対して，
  \[
    f_1(\bm{x} + \bm{y})
    = \begin{pmatrix} x_1 + y_1 \\ -(x_2 + y_2) \end{pmatrix}
    = \begin{pmatrix} x_1 \\ -x_2 \end{pmatrix} + \begin{pmatrix} y_1 \\ -y_2 \end{pmatrix}
    = f_1(\bm{x}) + f_1(\bm{y}),
  \]
  \[
    f_1(\alpha\bm{x})
    = \begin{pmatrix} \alpha x_1 \\ -\alpha x_2 \end{pmatrix}
    = \alpha \begin{pmatrix} x_1 \\ -x_2 \end{pmatrix}
    = \alpha f_1(\bm{x})
  \]
  が成り立つ．$f_2$，$f_3$についても同様である．
\end{example}

\begin{example}{原点のまわりの回転}{原点のまわりの回転}
  $K = \mathbb{R}$とし，$\theta$を定数とする．
  座標平面上の点を原点のまわりに角$\theta$だけ回転させる対応は，
  \[
    f\begin{pmatrix} x \\ y \end{pmatrix}
    = \begin{pmatrix} x\cos\theta - y\sin\theta \\ x\sin\theta + y\cos\theta \end{pmatrix}
  \]
  と表される\footnote{点$(x, y)$を極形式で$x = r\cos\varphi$，$y = r\sin\varphi$とかけば，回転後の点は$(r\cos(\varphi + \theta),\ r\sin(\varphi + \theta))$である．三角関数の加法定理を用いて展開すれば，上の式が得られる．}．
  この$f$は$\mathbb{R}^2$の線型変換である．実際，
  \[
    f\left( \begin{pmatrix} x_1 \\ x_2 \end{pmatrix} + \begin{pmatrix} y_1 \\ y_2 \end{pmatrix} \right)
    = \begin{pmatrix} (x_1 + y_1)\cos\theta - (x_2 + y_2)\sin\theta \\ (x_1 + y_1)\sin\theta + (x_2 + y_2)\cos\theta \end{pmatrix}
    = f\begin{pmatrix} x_1 \\ x_2 \end{pmatrix} + f\begin{pmatrix} y_1 \\ y_2 \end{pmatrix}
  \]
  であり，また$\alpha \in \mathbb{R}$に対して
  \[
    f\left( \alpha \begin{pmatrix} x_1 \\ x_2 \end{pmatrix} \right)
    = \begin{pmatrix} \alpha x_1 \cos\theta - \alpha x_2 \sin\theta \\ \alpha x_1 \sin\theta + \alpha x_2 \cos\theta \end{pmatrix}
    = \alpha f\begin{pmatrix} x_1 \\ x_2 \end{pmatrix}
  \]
  が成り立つからである．

  なお，この事実は幾何的にも納得できる．
  和$\bm{x} + \bm{y}$は平行四辺形の対角線として得られるが，平行四辺形全体を回転させても平行四辺形のままであり，
  対角線は対角線にうつる．これが$f(\bm{x} + \bm{y}) = f(\bm{x}) + f(\bm{y})$の意味するところである．
\end{example}

高校で学ぶベクトルの内積も，線型写像の例を与える．

\begin{example}{内積が定める線型写像}{内積が定める線型写像}
  $K = \mathbb{R}$とし，$\bm{a} = \begin{pmatrix} a_1 \\ a_2 \end{pmatrix} \in \mathbb{R}^2$を定ベクトルとする．
  写像$f \colon \mathbb{R}^2 \to \mathbb{R}$を，内積を用いて
  \[
    f(\bm{x}) = \bm{a} \cdot \bm{x} = a_1 x_1 + a_2 x_2
  \]
  で定めると，$f$は線型写像である．
  これは，高校で学ぶ内積の性質
  \[
    \bm{a} \cdot (\bm{x} + \bm{y}) = \bm{a} \cdot \bm{x} + \bm{a} \cdot \bm{y}, \qquad
    \bm{a} \cdot (\alpha\bm{x}) = \alpha (\bm{a} \cdot \bm{x})
  \]
  そのものにほかならない．

  さらに，$\bm{u}$を大きさ$1$のベクトルとするとき，$\bm{x}$から$\bm{u}$の定める直線への正射影は
  \[
    g(\bm{x}) = (\bm{u} \cdot \bm{x})\,\bm{u}
  \]
  で与えられる．内積の上の性質より，
  \[
    g(\bm{x} + \bm{y}) = \bigl(\bm{u} \cdot (\bm{x} + \bm{y})\bigr)\bm{u}
    = (\bm{u} \cdot \bm{x})\bm{u} + (\bm{u} \cdot \bm{y})\bm{u} = g(\bm{x}) + g(\bm{y})
  \]
  となり，スカラー倍についても同様であるから，$g$は$\mathbb{R}^2$の線型変換である．
  $\bm{u} = \bm{e}_1$とすれば，\exref{対称移動と正射影}の$f_3$が得られる．
\end{example}

線型性の定義から直ちに従う，簡単だが有用な性質を注意しておこう．

\begin{prop}{線型写像は原点を原点にうつす}{線型写像と零ベクトル}
  $f \colon K^n \to K^m$が線型写像ならば，
  \[
    f(\bm{0}) = \bm{0}
  \]
  が成り立つ．
\end{prop}

\begin{leftbar}
\begin{proof}
  \deref{線型写像}の第$2$の条件で$\alpha = 0$，$\bm{x} = \bm{0}$とすれば，
  \[
    f(\bm{0}) = f(0 \cdot \bm{0}) = 0 \cdot f(\bm{0}) = \bm{0}
  \]
  が成り立つ．これが証明すべきことであった．
\end{proof}
\end{leftbar}

\prref{線型写像と零ベクトル}の対偶をとれば，「$f(\bm{0}) \ne \bm{0}$ならば$f$は線型写像でない」となる．
これは，与えられた写像が線型でないことを示す際の，もっとも手軽な判定法である．

\begin{example}{線型写像でない写像}{線型写像でない写像}
  以下，$K = \mathbb{R}$とする．
  \begin{enumerate}[label=(\arabic*),leftmargin=3.2em]
    \item （平行移動）$\bm{b} \ne \bm{0}$を定ベクトルとし，$f \colon \mathbb{R}^2 \to \mathbb{R}^2$を$f(\bm{x}) = \bm{x} + \bm{b}$で定めると，
      $f(\bm{0}) = \bm{b} \ne \bm{0}$であるから，\prref{線型写像と零ベクトル}より$f$は線型写像でない．
    \item （$2$乗）$f \colon \mathbb{R} \to \mathbb{R}$を$f(x) = x^2$で定めると，
      \[
        f(1 + 1) = 4, \quad f(1) + f(1) = 2
      \]
      であり，$f(x + y) = f(x) + f(y)$が成り立たない．よって$f$は線型写像でない．
    \item （絶対値）$f \colon \mathbb{R} \to \mathbb{R}$を$f(x) = |x|$で定めると，$f(0) = 0$は成り立つが，
      \[
        f\bigl((-1) \cdot 1\bigr) = 1, \quad (-1) \cdot f(1) = -1
      \]
      であり，$f(\alpha x) = \alpha f(x)$が成り立たない．よって$f$は線型写像でない．
    \item （成分の積）$f \colon \mathbb{R}^2 \to \mathbb{R}$を$f\begin{pmatrix} x \\ y \end{pmatrix} = xy$で定めると，$f(\bm{0}) = 0$であるが，
      $\bm{x} = \begin{pmatrix} 1 \\ 1 \end{pmatrix}$に対して
      \[
        f(2\bm{x}) = 2 \cdot 2 = 4, \quad 2 f(\bm{x}) = 2 \cdot 1 = 2
      \]
      となり，やはり線型写像でない．
  \end{enumerate}
  (3)，(4)からわかるように，$f(\bm{0}) = \bm{0}$は線型写像であるための必要条件にすぎず，十分条件ではない．
\end{example}

\begin{mycolumn}
  高校数学では$y = ax + b$を\textbf{一次関数}とよぶ．
  しかし，$b \ne 0$のとき，この対応$f(x) = ax + b$は\prref{線型写像と零ベクトル}より線型写像ではない．
  「一次」という言葉と「線型」という言葉は，ここでずれるのである．

  グラフでいえば，線型写像$\mathbb{R} \to \mathbb{R}$に対応するのは「原点を通る直線」だけであり，
  切片$b$の分だけ平行移動された直線は線型写像を与えない．
  一般に，線型写像$f$と定ベクトル$\bm{b}$を用いて
  \[
    g(\bm{x}) = f(\bm{x}) + \bm{b}
  \]
  と表される写像を\textbf{アフィン写像}\index{あふぃんしゃぞう@アフィン写像}という．
  一次関数は，線型写像そのものではなく，アフィン写像の例なのである．

  なお，$b \ne 0$のときの$f(x) = ax + b$も，$1$次元分の「下駄」を履かせて
  \[
    \begin{pmatrix} x \\ 1 \end{pmatrix} \longmapsto \begin{pmatrix} ax + b \\ 1 \end{pmatrix}
  \]
  という$\mathbb{R}^2$の線型変換とみなすことができる．
  コンピュータグラフィックスで平行移動を行列で扱う際には，この考え方（同次座標）が用いられる．
\end{mycolumn}

幾何以外にも，高校数学のなかには線型写像が数多く隠れている．

\begin{example}{合計と平均}{合計と平均}
  $n$個のデータ$x_1, x_2, \dots, x_n$を$\bm{x} \in \mathbb{R}^n$とみなし，写像$S, M \colon \mathbb{R}^n \to \mathbb{R}$を
  \[
    S(\bm{x}) = x_1 + x_2 + \cdots + x_n, \qquad
    M(\bm{x}) = \frac{x_1 + x_2 + \cdots + x_n}{n}
  \]
  で定める．すなわち$S$は合計，$M$は平均である．
  このとき$S$，$M$はいずれも線型写像である．
  「$2$組のデータを足し合わせたものの平均は，それぞれの平均の和に等しい」「すべてのデータを$\alpha$倍すれば平均も$\alpha$倍になる」という，
  データの分析でおなじみの性質が，そのまま線型性である．

  一方，分散
  \[
    V(\bm{x}) = \frac{1}{n}\sum_{i=1}^{n} \bigl(x_i - M(\bm{x})\bigr)^2
  \]
  は線型写像ではない．実際，$V(\alpha\bm{x}) = \alpha^2 V(\bm{x})$であって$\alpha V(\bm{x})$ではない．
\end{example}

\begin{example}{多項式の微分と定積分}{多項式の微分と定積分}
  $2$次以下の多項式$a_0 + a_1 x + a_2 x^2$を，その係数を並べたベクトル
  \[
    \begin{pmatrix} a_0 \\ a_1 \\ a_2 \end{pmatrix} \in \mathbb{R}^3
  \]
  と同一視する．同様に，$1$次以下の多項式$b_0 + b_1 x$を$\begin{pmatrix} b_0 \\ b_1 \end{pmatrix} \in \mathbb{R}^2$と同一視する．

  多項式を微分する操作は，
  \[
    (a_0 + a_1 x + a_2 x^2)' = a_1 + 2a_2 x
  \]
  であるから，この同一視のもとで写像
  \[
    D \colon \mathbb{R}^3 \to \mathbb{R}^2, \quad
    D\begin{pmatrix} a_0 \\ a_1 \\ a_2 \end{pmatrix} = \begin{pmatrix} a_1 \\ 2a_2 \end{pmatrix}
  \]
  を定める．$D$は線型写像であり，これは高校で学ぶ微分の性質
  \[
    (f + g)' = f' + g', \qquad (\alpha f)' = \alpha f'
  \]
  にほかならない．

  同様に，定積分
  \[
    \int_0^1 (a_0 + a_1 x + a_2 x^2)\,dx = a_0 + \frac{a_1}{2} + \frac{a_2}{3}
  \]
  を対応させる写像$I \colon \mathbb{R}^3 \to \mathbb{R}$も線型写像である．
  これは，積分の線型性
  \[
    \int_0^1 \bigl(f(x) + g(x)\bigr) dx = \int_0^1 f(x)\,dx + \int_0^1 g(x)\,dx, \qquad
    \int_0^1 \alpha f(x)\,dx = \alpha \int_0^1 f(x)\,dx
  \]
  の言い換えである．
\end{example}

数列もまた，線型写像の格好の題材である．
以下では，有限個の項からなる数列$a_1, a_2, \dots, a_n$を，そのままベクトル
\[
  \bm{a} = \begin{pmatrix} a_1 \\ a_2 \\ \vdots \\ a_n \end{pmatrix} \in K^n
\]
とみなす．

\begin{example}{階差数列をとる写像}{階差数列をとる写像}
  写像$\Delta \colon K^n \to K^{n-1}$を，階差数列を対応させる写像
  \[
    \Delta \bm{a} =
    \begin{pmatrix}
      a_2 - a_1 \\ a_3 - a_2 \\ \vdots \\ a_n - a_{n-1}
    \end{pmatrix}
  \]
  で定める．このとき$\Delta$は線型写像である．
  実際，第$i$成分だけを見れば
  \[
    (a_{i+1} + b_{i+1}) - (a_i + b_i) = (a_{i+1} - a_i) + (b_{i+1} - b_i), \qquad
    \alpha a_{i+1} - \alpha a_i = \alpha (a_{i+1} - a_i)
  \]
  であり，これがすべての$i$について成り立つからである．

  「$2$つの数列を足してから階差をとっても，それぞれ階差をとってから足しても同じ」という，
  高校で意識せずに使っている事実が，この写像の線型性にほかならない．
\end{example}

\begin{example}{部分和をとる写像}{部分和をとる写像}
  写像$\Sigma \colon K^n \to K^n$を，初項から第$k$項までの和
  \[
    S_k = a_1 + a_2 + \cdots + a_k \quad (1 \le k \le n)
  \]
  を並べたベクトルを対応させる写像
  \[
    \Sigma \bm{a} = \begin{pmatrix} S_1 \\ S_2 \\ \vdots \\ S_n \end{pmatrix}
  \]
  で定める．各$S_k$は$a_1, \dots, a_n$の線型結合であるから，$\Sigma$は線型変換である．
  これは，和の記号についての性質
  \[
    \sum_{i=1}^{k} (a_i + b_i) = \sum_{i=1}^{k} a_i + \sum_{i=1}^{k} b_i, \qquad
    \sum_{i=1}^{k} \alpha a_i = \alpha \sum_{i=1}^{k} a_i
  \]
  の言い換えである．

  なお，$\Delta$が微分の，$\Sigma$が積分の離散版にあたることに注意しよう．
  高校で学ぶ関係式$a_n = S_n - S_{n-1}$は，「部分和をとってから階差をとると元に戻る」ことを述べており，
  微分積分学の基本定理の数列版とみることができる．
\end{example}

\begin{example}{等差数列を作る写像}{等差数列を作る写像}
  初項$a$と公差$d$の組$\begin{pmatrix} a \\ d \end{pmatrix} \in K^2$に対して，
  等差数列の第$n$項までを並べたベクトルを対応させる写像$\varphi \colon K^2 \to K^n$を
  \[
    \varphi \begin{pmatrix} a \\ d \end{pmatrix}
    = \begin{pmatrix} a \\ a + d \\ a + 2d \\ \vdots \\ a + (n-1)d \end{pmatrix}
  \]
  で定める．各成分は$a$と$d$の線型結合であるから，$\varphi$は線型写像である．

  すなわち，「初項どうし・公差どうしを足した等差数列は，もとの$2$つの等差数列を項ごとに足したものに等しい」ということである．

  一方，初項$a$と公比$r$の組から等比数列を作る写像
  \[
    \psi \begin{pmatrix} a \\ r \end{pmatrix} = \begin{pmatrix} a \\ ar \\ ar^2 \end{pmatrix}
  \]
  は線型写像ではない．実際，$\bm{x} = \begin{pmatrix} 1 \\ 2 \end{pmatrix}$に対して
  \[
    \psi(2\bm{x}) = \begin{pmatrix} 2 \\ 8 \\ 32 \end{pmatrix}, \qquad
    2\psi(\bm{x}) = \begin{pmatrix} 2 \\ 4 \\ 8 \end{pmatrix}
  \]
  であり，一致しない．
  公比が数列の「作られ方」そのものに関わるため，公比を変数とみるかぎり線型性は失われるのである．
  ただし，公比$r$を固定して初項だけを動かす写像$a \mapsto (a, ar, ar^2)$は線型写像である．
\end{example}

\begin{example}{数列に関する線型でない操作}{数列と線型でない操作}
  次の$2$つは，数列に対するごく自然な操作であるが，線型写像ではない．
  \begin{enumerate}[label=(\arabic*),leftmargin=3.2em]
    \item 各項を$2$乗する写像$\bm{a} \mapsto (a_1^2, a_2^2, \dots, a_n^2)$．
      $\alpha$倍すると各成分は$\alpha^2$倍になるので，$f(\alpha\bm{a}) = \alpha f(\bm{a})$が成り立たない．
    \item 隣り合う項の比$\bm{a} \mapsto (a_2/a_1, \dots, a_n/a_{n-1})$．
      そもそも$a_i = 0$のとき定義されず，$K^n$全体で定義された写像にすらならない．
  \end{enumerate}
  等差数列（差が一定）が線型な世界となじみ，等比数列（比が一定）がなじまないのは，このあたりの事情による．
\end{example}

\begin{mycolumn}
  \exref{多項式の微分と定積分}では多項式を係数の並びとみなし，本項の数列の例では数列を項の並びとみなすことで，それぞれ数ベクトルに翻訳した．
  同様の翻訳は，関数などさまざまな対象に対して可能である．
  第3章では，このような「翻訳」を経由せずに，多項式全体や関数全体をそのまま扱える枠組みとして\textbf{線型空間}を導入する．
  そこでは，微分や積分は線型空間のあいだの線型写像として，改めて位置づけられることになる．
\end{mycolumn}

次に，線型写像から新しい線型写像をつくる操作として，和・スカラー倍・合成を考える．

\begin{definition}{線型写像の和とスカラー倍}{線型写像の和とスカラー倍}
  $f, g \colon K^n \to K^m$を線型写像とする．
  このとき，
  \[
    h(\bm{x}) = f(\bm{x}) + g(\bm{x})
  \]
  で定まる写像$h \colon K^n \to K^m$を$f$と$g$の\textbf{和}\index{せんけいしゃぞうのわ@線型写像の和}とよび，$f + g$と表す．

  また，$\alpha \in K$と線型写像$f \colon K^n \to K^m$に対して，
  \[
    h(\bm{x}) = \alpha \cdot f(\bm{x})
  \]
  で定まる写像$h \colon K^n \to K^m$を$f$の\textbf{$\alpha$倍}\index{せんけいしゃぞうのすからーばい@線型写像のスカラー倍}とよび，$\alpha f$と表す．
\end{definition}

\begin{prop}{線型写像の和とスカラー倍の線型性}{線型写像の和とスカラー倍の線型性}
  $f, g \colon K^n \to K^m$が線型写像ならば，$f + g \colon K^n \to K^m$も線型写像である．
  また，$\alpha \in K$に対して，$\alpha f \colon K^n \to K^m$も線型写像である．
\end{prop}

\begin{leftbar}
\begin{proof}
  前半の主張を示す．任意の$\bm{x}, \bm{y} \in K^n$に対して，
  \begin{align}
    (f + g)(\bm{x} + \bm{y}) &= f(\bm{x} + \bm{y}) + g(\bm{x} + \bm{y}) \tag{$f+g$の定義} \\
    &= f(\bm{x}) + f(\bm{y}) + g(\bm{x}) + g(\bm{y}) \tag{$f,g$の線型性} \\
    &= \bigl(f(\bm{x}) + g(\bm{x})\bigr) + \bigl(f(\bm{y}) + g(\bm{y})\bigr) \nonumber \\
    &= (f + g)(\bm{x}) + (f + g)(\bm{y}) \tag{$f+g$の定義}
  \end{align}
  が成り立ち，さらに任意の$\alpha \in K$に対して，
  \begin{align}
    (f + g)(\alpha\bm{x}) &= f(\alpha\bm{x}) + g(\alpha\bm{x}) \tag{$f+g$の定義} \\
    &= \alpha f(\bm{x}) + \alpha g(\bm{x}) \tag{$f,g$の線型性} \\
    &= \alpha\bigl(f(\bm{x}) + g(\bm{x})\bigr) \nonumber \\
    &= \alpha(f + g)(\bm{x}) \tag{$f+g$の定義}
  \end{align}
  が成り立つ．よって$f + g$は\deref{線型写像}の条件を満たし，線型写像である．

  後半の主張も同様にして証明される．
\end{proof}
\end{leftbar}

\begin{prop}{線型写像の合成}{線型写像の合成}
  $f \colon K^n \to K^m$，$g \colon K^m \to K^l$を線型写像とする．
  このとき，合成写像$g \circ f \colon K^n \to K^l$も線型写像である．
\end{prop}

\begin{leftbar}
\begin{proof}
  任意の$\bm{x}, \bm{y} \in K^n$に対して，
  \begin{align}
    (g \circ f)(\bm{x} + \bm{y}) &= g\bigl(f(\bm{x} + \bm{y})\bigr) \nonumber \\
    &= g\bigl(f(\bm{x}) + f(\bm{y})\bigr) \tag{$f$の線型性} \\
    &= g\bigl(f(\bm{x})\bigr) + g\bigl(f(\bm{y})\bigr) \tag{$g$の線型性} \\
    &= (g \circ f)(\bm{x}) + (g \circ f)(\bm{y}) \nonumber
  \end{align}
  が成り立ち，さらに任意の$\alpha \in K$に対して，
  \begin{align}
    (g \circ f)(\alpha\bm{x}) &= g\bigl(f(\alpha\bm{x})\bigr) \nonumber \\
    &= g\bigl(\alpha f(\bm{x})\bigr) \tag{$f$の線型性} \\
    &= \alpha g\bigl(f(\bm{x})\bigr) \tag{$g$の線型性} \\
    &= \alpha(g \circ f)(\bm{x}) \nonumber
  \end{align}
  が成り立つ．これが証明すべきことであった．
\end{proof}
\end{leftbar}

\phantomsection
\subsection{線型写像と行列}

本節では，線型写像と行列の関係を調べる．
結論からいえば，$K^n$から$K^m$への線型写像は$m \times n$行列とちょうど一対一に対応する．

以下では，第$j$成分が$1$で他の成分がすべて$0$であるベクトル
\[
  \bm{e}_j =
  \begin{pmatrix}
    0 \\ \vdots \\ 1 \\ \vdots \\ 0
  \end{pmatrix}
  \in K^n
  \quad (j = 1, 2, \dots, n)
\]
を\textbf{標準基底ベクトル}\index{ひょうじゅんきていべくとる@標準基底ベクトル}とよぶ\footnote{「基底」という概念そのものは，第3章で正式に定義される（\deref{基底}）．そこで見るように，$(\bm{e}_1, \bm{e}_2, \dots, \bm{e}_n)$は$K^n$の基底のもっとも代表的な例であり，「標準基底ベクトル」という名称もそれに由来する．現段階では，「第$j$成分だけが$1$で他の成分がすべて$0$のベクトル」を表す名前だと思っておけばよい．}．

\begin{prop}{行列の定める線型写像}{行列の定める線型写像}
  $A$を$m \times n$行列とする．
  $\bm{x} \in K^n$に対して$A\bm{x} \in K^m$を対応させる写像を$T_A$とかく．
  このとき，$T_A$は$K^n$から$K^m$への線型写像である．
\end{prop}

\begin{leftbar}
\begin{proof}
  任意の$\bm{x}, \bm{y} \in K^n$と$\alpha \in K$に対して，
  \begin{align}
    T_A(\bm{x} + \bm{y}) &= A(\bm{x} + \bm{y}) \tag{$T_A$の定義} \\
    &= A\bm{x} + A\bm{y} \tag{行列の積の分配律} \\
    &= T_A(\bm{x}) + T_A(\bm{y}) \nonumber
  \end{align}
  が成り立ち，さらに，
  \begin{align}
    T_A(\alpha\bm{x}) &= A(\alpha\bm{x}) \tag{$T_A$の定義} \\
    &= \alpha(A\bm{x}) \tag{行列の積の性質} \\
    &= \alpha T_A(\bm{x}) \nonumber
  \end{align}
  が成り立つ．
  よって$T_A$は線型写像である．
\end{proof}
\end{leftbar}

\prref{行列の定める線型写像}は「行列が与えられれば線型写像が定まる」ことを主張している．
実は，この逆も成り立つ．すなわち，任意の線型写像はある行列から定まる線型写像として表され，しかもそのような行列はただひとつしか存在しない．

\begin{theorem}{線型写像の行列表示}{線型写像の行列表示}
  $T \colon K^n \to K^m$を線型写像とする．
  このとき，ある$m \times n$行列$A$が存在して，任意の$\bm{x} \in K^n$に対して
  \[
    T(\bm{x}) = A\bm{x}
  \]
  が成り立つ．さらに，このような$A$はただひとつしか存在しない．
\end{theorem}

\begin{leftbar}
\begin{proof}
  まず，存在を示す．
  \[
    \bm{a}_j = T(\bm{e}_j) \quad (1 \le j \le n)
  \]
  とおき，これらを列として並べた$m \times n$行列
  \[
    A = (\bm{a}_1, \bm{a}_2, \dots, \bm{a}_n)
  \]
  を考える．$A$の定める線型写像を$T_A \colon K^n \to K^m$とすると，$1 \le j \le n$の範囲で
  \begin{equation}
    T_A(\bm{e}_j) = A\bm{e}_j = \bm{a}_j = T(\bm{e}_j) \label{eq:行列表示の基底での一致}
  \end{equation}
  が成り立つ．ここで，任意の$\bm{x} = \begin{pmatrix} x_1 \\ x_2 \\ \vdots \\ x_n \end{pmatrix} \in K^n$は
  \[
    \bm{x} = \sum_{j=1}^{n} x_j \bm{e}_j
  \]
  とかけるので，$T_A$と$T$の線型性および\eqref{eq:行列表示の基底での一致}より，
  \begin{align*}
    T_A(\bm{x}) &= T_A\left(\sum_{j=1}^{n} x_j \bm{e}_j\right)
    = \sum_{j=1}^{n} x_j T_A(\bm{e}_j)
    = \sum_{j=1}^{n} x_j T(\bm{e}_j)
    = T\left(\sum_{j=1}^{n} x_j \bm{e}_j\right)
    = T(\bm{x})
  \end{align*}
  が成り立つ．よって$T = T_A$であり，$A$の存在が示された．

  次に，一意性を示す．$A$，$B$がともに$m \times n$行列であって$T_A = T_B$を満たすとし，
  \[
    A = (\bm{a}_1, \bm{a}_2, \dots, \bm{a}_n), \quad
    B = (\bm{b}_1, \bm{b}_2, \dots, \bm{b}_n)
  \]
  とかく．このとき，$1 \le j \le n$の範囲で
  \begin{align}
    \bm{a}_j &= A\bm{e}_j \nonumber \\
    &= T_A(\bm{e}_j) \tag{$T_A$の定義} \\
    &= T_B(\bm{e}_j) \tag{$T_A = T_B$} \\
    &= B\bm{e}_j = \bm{b}_j \nonumber
  \end{align}
  が成り立つ．すべての列が一致するので$A = B$であり，一意性が示された．
  これが証明すべきことであった．
\end{proof}
\end{leftbar}

まとめると，行列$A$が与えられたとき，それに対応する写像$T_A$は線型写像であり，逆に，$n$次元のベクトルを$m$次元のベクトルにうつす任意の線型写像は，必ずただひとつの行列によって表される．
この意味で，「$K^n$から$K^m$への線型写像」と「$m \times n$行列」は同一視できる．

\begin{mycolumn}
  \thref{線型写像の行列表示}の証明の核心は，「線型写像は標準基底ベクトルの行き先だけで完全に決まる」という点にある．
  実際，任意のベクトルは標準基底ベクトルの線型結合（次節で定義する）として表せるので，線型性を用いれば，$T(\bm{e}_1), T(\bm{e}_2), \dots, T(\bm{e}_n)$の値から$T(\bm{x})$の値がすべて復元できる．
  この「基底での値がすべてを決める」という考え方は，線型代数の随所で現れる重要な視点である．
\end{mycolumn}

この同一視のもとで，行列の積は写像の合成に対応する．

\begin{prop}{行列の積と写像の合成}{行列の積と写像の合成}
  $A$を$l \times m$行列，$B$を$m \times n$行列とする．このとき，
  \[
    T_A \circ T_B = T_{AB}
  \]
  が成り立つ．
\end{prop}

\begin{leftbar}
\begin{proof}
  任意の$\bm{x} \in K^n$に対して，
  \begin{align}
    (T_A \circ T_B)(\bm{x}) &= T_A\bigl(T_B(\bm{x})\bigr) \nonumber \\
    &= A(B\bm{x}) \tag{$T_A$，$T_B$の定義} \\
    &= (AB)\bm{x} \tag{\prref{行列の乗法結合律}} \\
    &= T_{AB}(\bm{x}) \nonumber
  \end{align}
  が成り立つ\footnote{$\bm{x} \in K^n$を$n \times 1$行列とみなせば，$A(B\bm{x}) = (AB)\bm{x}$は行列の乗法結合律そのものである．}．
  これが証明すべきことであった．
\end{proof}
\end{leftbar}

\prref{行列の積と写像の合成}によって，一見複雑に見えた行列の積の定義（\deref{行列の積}）が，「写像の合成に対応するように定められたもの」として自然に理解できる．
行列の積が交換律を満たさないのは，写像の合成が交換律を満たさないことの反映なのである．

\thref{線型写像の行列表示}の証明は，行列$A$の第$j$列が$T(\bm{e}_j)$であることを教えてくれる．
これを用いれば，前節で見た線型変換を表す行列は，標準基底ベクトルの行き先を調べるだけで求められる．

\begin{example}{図形の移動を表す行列}{図形の移動を表す行列}
  $K = \mathbb{R}$とする．
  \exref{対称移動と正射影}の$f_1$（$x$軸に関する対称移動）については，
  \[
    f_1(\bm{e}_1) = \begin{pmatrix} 1 \\ 0 \end{pmatrix}, \quad
    f_1(\bm{e}_2) = \begin{pmatrix} 0 \\ -1 \end{pmatrix}
  \]
  であるから，これらを列として並べて
  \[
    f_1(\bm{x}) = \begin{pmatrix} 1 & 0 \\ 0 & -1 \end{pmatrix} \bm{x}
  \]
  となる．同様にして，
  \[
    f_2(\bm{x}) = \begin{pmatrix} 0 & 1 \\ 1 & 0 \end{pmatrix} \bm{x}
    \quad \text{（直線$y = x$に関する対称移動）}, \qquad
    f_3(\bm{x}) = \begin{pmatrix} 1 & 0 \\ 0 & 0 \end{pmatrix} \bm{x}
    \quad \text{（$x$軸への正射影）}
  \]
  が得られる．
  また，\exref{原点のまわりの回転}の回転については，
  \[
    f(\bm{e}_1) = \begin{pmatrix} \cos\theta \\ \sin\theta \end{pmatrix}, \quad
    f(\bm{e}_2) = \begin{pmatrix} -\sin\theta \\ \cos\theta \end{pmatrix}
  \]
  であるから，
  \[
    R(\theta) = \begin{pmatrix} \cos\theta & -\sin\theta \\ \sin\theta & \cos\theta \end{pmatrix}
  \]
  とおけば$f(\bm{x}) = R(\theta)\bm{x}$である．この$R(\theta)$を\textbf{回転行列}\index{かいてんぎょうれつ@回転行列}という．
\end{example}

\begin{example}{回転行列と加法定理}{回転行列と加法定理}
  角$\varphi$の回転をしてから角$\theta$の回転をすれば，全体としては角$\theta + \varphi$の回転になる．
  これを\prref{行列の積と写像の合成}によって行列の言葉に翻訳すると，
  \[
    R(\theta) R(\varphi) = R(\theta + \varphi)
  \]
  となる．左辺を\deref{行列の積}に従って計算すれば，
  \[
    R(\theta)R(\varphi) =
    \begin{pmatrix}
      \cos\theta\cos\varphi - \sin\theta\sin\varphi & -(\cos\theta\sin\varphi + \sin\theta\cos\varphi) \\
      \sin\theta\cos\varphi + \cos\theta\sin\varphi & \cos\theta\cos\varphi - \sin\theta\sin\varphi
    \end{pmatrix}
  \]
  であり，右辺は
  \[
    R(\theta + \varphi) =
    \begin{pmatrix}
      \cos(\theta + \varphi) & -\sin(\theta + \varphi) \\
      \sin(\theta + \varphi) & \cos(\theta + \varphi)
    \end{pmatrix}
  \]
  である．両辺の成分を比較すれば，三角関数の加法定理
  \[
    \cos(\theta + \varphi) = \cos\theta\cos\varphi - \sin\theta\sin\varphi, \qquad
    \sin(\theta + \varphi) = \sin\theta\cos\varphi + \cos\theta\sin\varphi
  \]
  が得られる．
  行列の積という「一見複雑な定義」が，写像の合成という自然な操作を正確に写し取っていることの，ひとつの現れである．
\end{example}

\begin{example}{回転行列の逆行列}{回転行列の逆行列}
  角$\theta$の回転を打ち消すのは角$-\theta$の回転であるから，
  \[
    R(\theta) R(-\theta) = R(0) = E_2
  \]
  が成り立つ．よって$R(\theta)$は正則であり，
  \[
    R(\theta)^{-1} = R(-\theta) = \begin{pmatrix} \cos\theta & \sin\theta \\ -\sin\theta & \cos\theta \end{pmatrix}
    = R(\theta)^\mathsf{T}
  \]
  である．
  一方，\exref{対称移動と正射影}の正射影$f_3$を表す行列$\begin{pmatrix} 1 & 0 \\ 0 & 0 \end{pmatrix}$は正則でない．
  実際，$\bm{e}_2 \ne \bm{0}$が$\bm{0}$にうつるため，$f_3$は逆をもちえない．
  「つぶれてしまう変換は元に戻せない」という直観が，逆行列の存在・非存在に対応しているのである．
\end{example}

\begin{mycolumn}
  複素数平面もまた，線型変換の例を与える．
  $\alpha = p + qi$を定数とし，複素数$z = x + yi$に$\alpha z$を対応させる写像を考えよう．
  \[
    \alpha z = (p + qi)(x + yi) = (px - qy) + (qx + py)i
  \]
  であるから，$z = x + yi$を$\begin{pmatrix} x \\ y \end{pmatrix} \in \mathbb{R}^2$と同一視すれば，この対応は
  \[
    \begin{pmatrix} x \\ y \end{pmatrix} \longmapsto
    \begin{pmatrix} p & -q \\ q & p \end{pmatrix}
    \begin{pmatrix} x \\ y \end{pmatrix}
  \]
  という$\mathbb{R}^2$の線型変換にほかならない．
  とくに$\alpha = \cos\theta + i\sin\theta$のとき，この行列は回転行列$R(\theta)$と一致する．
  「複素数を極形式でかけ算すると偏角が足される」という高校で学ぶ事実と，
  \exref{回転行列と加法定理}の$R(\theta)R(\varphi) = R(\theta + \varphi)$とが，同じ内容を述べていることがわかる．
\end{mycolumn}

数列に関する例も，行列の言葉で見直してみよう．

\begin{example}{階差行列と部分和行列}{階差行列と部分和行列}
  \exref{部分和をとる写像}の$\Sigma \colon K^n \to K^n$を表す行列を求めよう．
  $\Sigma \bm{e}_j$は「第$j$項だけが$1$で他が$0$の数列」の部分和であるから，
  第$j$成分以降がすべて$1$，それより前が$0$のベクトルである．
  これらを列として並べると，$n = 4$のとき
  \[
    L =
    \begin{pmatrix}
      1 & 0 & 0 & 0 \\
      1 & 1 & 0 & 0 \\
      1 & 1 & 1 & 0 \\
      1 & 1 & 1 & 1
    \end{pmatrix}
  \]
  となる（対角成分より下がすべて$1$の行列である）．

  一方，\exref{階差数列をとる写像}の$\Delta$を少し変形し，「第$1$項はそのまま残し，第$2$項以降は階差をとる」写像
  \[
    \widetilde{\Delta} \bm{a} =
    \begin{pmatrix} a_1 \\ a_2 - a_1 \\ a_3 - a_2 \\ a_4 - a_3 \end{pmatrix}
  \]
  を考えると，これは$K^4$の線型変換であり，その行列は
  \[
    D =
    \begin{pmatrix}
      1 & 0 & 0 & 0 \\
      -1 & 1 & 0 & 0 \\
      0 & -1 & 1 & 0 \\
      0 & 0 & -1 & 1
    \end{pmatrix}
  \]
  である．実際に積を計算すれば
  \[
    DL = LD = E_4
  \]
  となり，$D = L^{-1}$であることがわかる．
  これは，「部分和をとってから階差をとると元の数列に戻る」という\exref{部分和をとる写像}の観察を，
  逆行列の言葉で述べ直したものである．
\end{example}

\begin{example}{漸化式が定める線型変換}{漸化式が定める線型変換}
  $p, q \in K$を定数とし，数列$\{a_n\}$が三項間漸化式
  \[
    a_{n+2} = p\,a_{n+1} + q\,a_n
  \]
  を満たすとする．
  連続する$2$項を並べたベクトル$\begin{pmatrix} a_n \\ a_{n+1} \end{pmatrix} \in K^2$に着目すると，
  漸化式は
  \[
    \begin{pmatrix} a_{n+1} \\ a_{n+2} \end{pmatrix}
    = \begin{pmatrix} 0 & 1 \\ q & p \end{pmatrix}
      \begin{pmatrix} a_{n} \\ a_{n+1} \end{pmatrix}
  \]
  と表される．
  すなわち，「$1$項だけ先へ進める」という操作が$K^2$の線型変換にほかならない．

  この行列を$A$とおくと，\prref{行列の積と写像の合成}をくり返し用いて
  \[
    \begin{pmatrix} a_{n+1} \\ a_{n+2} \end{pmatrix}
    = A^{n} \begin{pmatrix} a_{1} \\ a_{2} \end{pmatrix}
  \]
  が得られる．
  たとえばフィボナッチ数列（$F_1 = F_2 = 1$，$F_{n+2} = F_{n+1} + F_n$）であれば，
  \[
    A = \begin{pmatrix} 0 & 1 \\ 1 & 1 \end{pmatrix}, \qquad
    \begin{pmatrix} F_{n+1} \\ F_{n+2} \end{pmatrix}
    = A^{n} \begin{pmatrix} 1 \\ 1 \end{pmatrix}
  \]
  である．漸化式を$n$回追いかける代わりに，行列の$n$乗を計算すればよいことになる．
\end{example}

\begin{mycolumn}
  \exref{漸化式が定める線型変換}によって，数列の一般項を求める問題は「行列$A$の$n$乗を求める問題」に翻訳された．
  では，$A^n$はどう計算すればよいのだろうか．

  鍵となるのは，$A\bm{x} = \lambda \bm{x}$を満たすベクトル$\bm{x} \ne \bm{0}$である．
  このような$\bm{x}$に対しては$A^n \bm{x} = \lambda^n \bm{x}$となり，べき乗が一気に見通しよくなる．
  この$\lambda$を$A$の固有値とよび，第4章で詳しく扱う．
  高校で三項間漸化式を解くときに現れる$2$次方程式$t^2 = pt + q$（特性方程式）は，
  実は$A$の固有値を求める方程式にほかならない．

  また，漸化式$a_{n+2} = p a_{n+1} + q a_n$を満たす数列全体の集合を考えると，
  これは項ごとの和とスカラー倍について閉じている．
  高校で学ぶ「$2$つの解の線型結合もまた解である」という事実は，この集合が線型空間をなすことを意味しており，
  さらにその次元が$2$である（初項と第$2$項で解が決まる）ことが，一般解が$2$つの解の線型結合でつくされる理由である．
  これらの概念は第3章で整備される．
\end{mycolumn}

\phantomsection
\subsection{線型写像の像と核}

線型写像$f \colon K^n \to K^m$が与えられると，「$f$の行き先として実際に現れるベクトル全体」と「$f$によって$\bm{0}$にうつされるベクトル全体」という二つの集合が自然に定まる．
本項では，この二つを\textbf{像}・\textbf{核}として定式化する．
これらは，第3章の次元公式（\thref{線型写像の次元公式}）をはじめ，以後の議論の随所で用いられる基本概念である．

その準備として，まず一般の写像に関する用語を整えておこう．

\begin{definition}{単射・全射・全単射}{単射と全射}
  $f \colon V \to W$を写像とする．
  \begin{enumerate}[label=(\roman*),leftmargin=3.2em]
    \item 任意の$x, y \in V$に対して「$f(x) = f(y)$ならば$x = y$」が成り立つとき，$f$は\textbf{単射}\index{たんしゃ@単射}であるという．
    \item 任意の$y \in W$に対して，$f(x) = y$となる$x \in V$が存在するとき，$f$は\textbf{全射}\index{ぜんしゃ@全射}であるという．
    \item $f$が単射かつ全射であるとき，$f$は\textbf{全単射}\index{ぜんたんしゃ@全単射}であるという．
  \end{enumerate}
\end{definition}

単射とは「異なる元は必ず異なる元にうつる」こと，全射とは「終域のどの元も行き先として現れる」ことである\footnote{全単射は，第2章で置換を定義する際にも用いられる．}．

\begin{definition}{像と核}{像と核}
  $f \colon K^n \to K^m$を線型写像とする．
  \[
    \Im f = \{ f(\bm{x}) \mid \bm{x} \in K^n \} \subset K^m
  \]
  を$f$の\textbf{像}\index{ぞう@像}\index{Im@$\Im$}（image）という．
  また，
  \[
    \Ker f = \{ \bm{x} \in K^n \mid f(\bm{x}) = \bm{0} \} \subset K^n
  \]
  を$f$の\textbf{核}\index{かく@核}\index{Ker@$\Ker$}（kernel）という．
\end{definition}

像が終域$K^m$の部分集合であるのに対し，核は定義域$K^n$の部分集合である．
両者の「住んでいる場所」の違いに注意してほしい\footnote{像は，線型とは限らない一般の写像$f \colon V \to W$に対しても同じ式で定義される．一方，核の定義には終域に$\bm{0}$という特別な元があることが必要であり，こちらは線型写像に固有の概念である．}．

\prref{線型写像と零ベクトル}より$f(\bm{0}) = \bm{0}$であるから，つねに$\bm{0} \in \Ker f$かつ$\bm{0} \in \Im f$であり，とくに像も核も空集合ではない．
また，定義から直ちに，$f$が全射であることと$\Im f = K^m$であることは同値である．

\begin{example}{像と核}{像と核の例}
  \begin{enumerate}[label=(\arabic*),leftmargin=3.2em]
    \item 恒等変換$\mathrm{id} \colon K^n \to K^n$（\exref{恒等変換と零写像}）については，
      \[
        \Im \mathrm{id} = K^n, \qquad \Ker \mathrm{id} = \{\bm{0}\}
      \]
      である．一方，零写像$f \colon K^n \to K^m$については，
      \[
        \Im f = \{\bm{0}\}, \qquad \Ker f = K^n
      \]
      である．
    \item $x$軸への正射影$f_3 \colon \mathbb{R}^2 \to \mathbb{R}^2$（\exref{対称移動と正射影}）については，
      \[
        \Im f_3 = \left\{ \begin{pmatrix} x \\ 0 \end{pmatrix} ~\middle|~ x \in \mathbb{R} \right\}, \qquad
        \Ker f_3 = \left\{ \begin{pmatrix} 0 \\ y \end{pmatrix} ~\middle|~ y \in \mathbb{R} \right\}
      \]
      である．すなわち，像は$x$軸，核は$y$軸である．
      平面全体が$x$軸へとつぶされるとき，「届く範囲」が像であり，「原点へつぶれる方向」が核なのである．
    \item 多項式の微分$D \colon \mathbb{R}^3 \to \mathbb{R}^2$（\exref{多項式の微分と定積分}）については，
      \[
        \Ker D = \left\{ \begin{pmatrix} a_0 \\ 0 \\ 0 \end{pmatrix} ~\middle|~ a_0 \in \mathbb{R} \right\}
      \]
      である．これは「（$2$次以下の）多項式のうち，微分して$0$になるのは定数にかぎる」という事実の言い換えである．
      また，$1$次以下の任意の多項式$b_0 + b_1 x$は，たとえば$b_0 x + \frac{b_1}{2} x^2$の微分として得られる（原始関数の存在）から，$\Im D = \mathbb{R}^2$である．
  \end{enumerate}
\end{example}

連立一次方程式のことばを使えば，行列の定める線型写像$T_A$の像と核は，すでによく知っている対象であることがわかる．
以下，右辺の定数項がすべて$0$である連立一次方程式$A\bm{x} = \bm{0}$を\textbf{斉次連立一次方程式}\index{せいじれんりついちじほうていしき@斉次連立一次方程式}（同次連立一次方程式）という．

\begin{prop}{連立一次方程式と像・核}{連立一次方程式と像と核}
  $A$を$m \times n$行列とし，$T_A \colon K^n \to K^m$を$A$の定める線型写像（\prref{行列の定める線型写像}）とする．このとき，次が成り立つ．
  \begin{enumerate}[label=(\roman*),leftmargin=3.2em]
    \item $\Ker T_A$は，斉次連立一次方程式$A\bm{x} = \bm{0}$の解全体の集合である．
    \item $\bm{b} \in K^m$に対して，$\bm{b} \in \Im T_A$であることと，連立一次方程式$A\bm{x} = \bm{b}$が解をもつことは同値である．
  \end{enumerate}
\end{prop}

\begin{leftbar}
\begin{proof}
  (i)について，$\bm{x} \in \Ker T_A$であることは，$T_A$の定義より$A\bm{x} = \bm{0}$と同値であり，これは$\bm{x}$が$A\bm{x} = \bm{0}$の解であることにほかならない．

  (ii)について，$\bm{b} \in \Im T_A$であることは，ある$\bm{x} \in K^n$が存在して$A\bm{x} = \bm{b}$となること，すなわち$A\bm{x} = \bm{b}$が解をもつことと同値である．
  いずれも定義の言い換えである．これが証明すべきことであった．
\end{proof}
\end{leftbar}

このように，核は「斉次連立一次方程式の解全体」を，像は「右辺として実現可能なベクトル全体」を，写像のことばで言い直したものである．

核のもっとも重要な役割のひとつは，線型写像の単射性の判定である．

\begin{prop}{単射性と核}{単射性と核}
  線型写像$f \colon K^n \to K^m$に対して，次の二つは同値である．
  \begin{enumerate}[label=(\roman*),leftmargin=3.2em]
    \item $f$は単射である．
    \item $\Ker f = \{\bm{0}\}$である．
  \end{enumerate}
\end{prop}

\begin{leftbar}
\begin{proof}
  (i)$\Rightarrow$(ii)：$\bm{0} \in \Ker f$はつねに成り立つので，逆の包含を示せばよい．
  $\bm{x} \in \Ker f$とすると，\prref{線型写像と零ベクトル}より
  \[
    f(\bm{x}) = \bm{0} = f(\bm{0})
  \]
  である．$f$は単射であるから$\bm{x} = \bm{0}$となり，$\Ker f = \{\bm{0}\}$である．

  (ii)$\Rightarrow$(i)：$f(\bm{x}) = f(\bm{y})$とする．$f$の線型性より
  \[
    f(\bm{x} - \bm{y}) = f\bigl(\bm{x} + (-1)\bm{y}\bigr) = f(\bm{x}) + (-1)f(\bm{y}) = f(\bm{x}) - f(\bm{y}) = \bm{0}
  \]
  であるから，$\bm{x} - \bm{y} \in \Ker f = \{\bm{0}\}$である．
  よって$\bm{x} - \bm{y} = \bm{0}$，すなわち$\bm{x} = \bm{y}$となり，$f$は単射である．
  これが証明すべきことであった．
\end{proof}
\end{leftbar}

一般の写像の単射性を確かめるには，すべての値の組$f(\bm{x})$，$f(\bm{y})$を比較しなければならない．
これに対して線型写像では，\prref{単射性と核}により「$\bm{0}$にうつるベクトルが$\bm{0}$だけかどうか」を調べれば十分である．
これが核という概念の威力である．

\exref{回転行列の逆行列}では，正射影$f_3$について「$\bm{e}_2 \ne \bm{0}$が$\bm{0}$にうつるため，$f_3$は逆をもちえない」と述べた．
これはまさに$\bm{e}_2 \in \Ker f_3$，すなわち$\Ker f_3 \ne \{\bm{0}\}$という主張であり，\prref{単射性と核}より$f_3$は単射でない．
「つぶれてしまう変換は元に戻せない」という直観は，核のことばで正確に表現されるのである．

\begin{mycolumn}
  像と核は，この後の章で中心的な役割を果たす．
  第3章では，$\Im f$と$\Ker f$が単なる部分集合ではなく，それぞれ$K^m$，$K^n$の\textbf{線型部分空間}であることを確かめる（\prref{像と核の線型部分空間性}）．
  さらに「次元」を導入すると，両者のあいだには
  \[
    \dim \Ker f + \dim \Im f = n
  \]
  という関係が成り立つ（線型写像の次元公式，\thref{線型写像の次元公式}）．
  定義域のもつ$n$次元分の情報のうち，核の分だけが$\bm{0}$につぶれ，残りが像として生き残る，という直観を正確に述べたのがこの公式である．
\end{mycolumn}

\phantomsection
\section{線型結合}

\phantomsection
\subsection{線型結合と線型関係}

本節では，複数のベクトルからつくられる基本的な演算である\textbf{線型結合}\index{せんけいけつごう@線型結合}を定義する\footnote{ここでは$K^m$のベクトルに対して定義するが，のちの章で導入する一般の線型空間の元に対しても，まったく同様に定義される．}．

\begin{definition}{線型結合・線型関係}{線型結合}
  $n$個のベクトル$\bm{x}_1, \bm{x}_2, \dots, \bm{x}_n \in K^m$を考える．
  \begin{equation}
    \lambda_1 \bm{x}_1 + \lambda_2 \bm{x}_2 + \cdots + \lambda_n \bm{x}_n
    \quad (\lambda_i \in K) \label{eq:線型結合の形}
  \end{equation}
  の形のベクトルを，$\bm{x}_1, \bm{x}_2, \dots, \bm{x}_n$の\textbf{線型結合}という．

  また，\eqref{eq:線型結合の形}に関連して，等式
  \[
    \lambda_1 \bm{x}_1 + \lambda_2 \bm{x}_2 + \cdots + \lambda_n \bm{x}_n = \bm{0}
  \]
  を$\bm{x}_1, \bm{x}_2, \dots, \bm{x}_n$の\textbf{線型関係}\index{せんけいかんけい@線型関係}という．
\end{definition}

任意のベクトル$\bm{x}_1, \bm{x}_2, \dots, \bm{x}_n$に対して，係数をすべて$0$とした
\[
  0\bm{x}_1 + 0\bm{x}_2 + \cdots + 0\bm{x}_n = \bm{0}
\]
はつねに成り立つ線型関係である．これを\textbf{自明な線型関係}\index{じめいなせんけいかんけい@自明な線型関係}とよぶ．

前節で導入した線型写像は，線型結合と相性がよい．
すなわち，線型写像は線型結合を保つ：

\begin{prop}{線型写像と線型結合}{線型写像と線型結合}
  $f \colon K^n \to K^m$を線型写像とする．
  このとき，任意の$\bm{x}_1, \bm{x}_2, \dots, \bm{x}_k \in K^n$と任意の$\lambda_1, \lambda_2, \dots, \lambda_k \in K$に対して，
  \[
    f(\lambda_1 \bm{x}_1 + \lambda_2 \bm{x}_2 + \cdots + \lambda_k \bm{x}_k)
    = \lambda_1 f(\bm{x}_1) + \lambda_2 f(\bm{x}_2) + \cdots + \lambda_k f(\bm{x}_k)
  \]
  が成り立つ．
\end{prop}

\begin{leftbar}
\begin{proof}
  $k$に関する帰納法で示す．
  $k = 1$のときは，$f$の線型性の第$2$の条件そのものである．

  $k - 1$のとき主張が成り立つと仮定する．このとき，
  \begin{align}
    f\left(\sum_{i=1}^{k} \lambda_i \bm{x}_i\right)
    &= f\left(\sum_{i=1}^{k-1} \lambda_i \bm{x}_i + \lambda_k \bm{x}_k\right) \nonumber \\
    &= f\left(\sum_{i=1}^{k-1} \lambda_i \bm{x}_i\right) + f(\lambda_k \bm{x}_k) \tag{$f$の線型性} \\
    &= \sum_{i=1}^{k-1} \lambda_i f(\bm{x}_i) + \lambda_k f(\bm{x}_k) \tag{帰納法の仮定と$f$の線型性} \\
    &= \sum_{i=1}^{k} \lambda_i f(\bm{x}_i) \nonumber
  \end{align}
  となり，$k$のときも主張が成り立つ．これが証明すべきことであった．
\end{proof}
\end{leftbar}

\begin{mycolumn}
  行列とベクトルの積は，線型結合のことばで見直すことができる．
  $m \times n$行列$A$の列ベクトルを$\bm{a}_1, \bm{a}_2, \dots, \bm{a}_n \in K^m$とすると，任意の$\bm{x} = \begin{pmatrix} x_1 \\ x_2 \\ \vdots \\ x_n \end{pmatrix} \in K^n$に対して，
  \[
    A\bm{x} = x_1 \bm{a}_1 + x_2 \bm{a}_2 + \cdots + x_n \bm{a}_n
  \]
  が成り立つ．実際，$\bm{x} = \sum_{j=1}^{n} x_j \bm{e}_j$と\prref{線型写像と線型結合}より，
  \[
    A\bm{x} = T_A\left(\sum_{j=1}^{n} x_j \bm{e}_j\right)
    = \sum_{j=1}^{n} x_j T_A(\bm{e}_j)
    = \sum_{j=1}^{n} x_j \bm{a}_j
  \]
  となる．
  すなわち，$A\bm{x}$は$A$の列ベクトルの線型結合であり，その係数は$\bm{x}$の成分にほかならない．
  この見方によれば，連立一次方程式$A\bm{x} = \bm{b}$が解を持つことは，$\bm{b}$が$A$の列ベクトルの線型結合として表せることと同値である．
  一方，\prref{連立一次方程式と像と核}(ii)より，$A\bm{x} = \bm{b}$が解を持つことは$\bm{b} \in \Im T_A$とも同値であった．
  あわせると，像$\Im T_A$とは「$A$の列ベクトルの線型結合全体の集合」にほかならないことがわかる．
\end{mycolumn}

\phantomsection
\subsection{線型独立と線型従属}

線型関係のうち，自明なものしか存在しないとき，ベクトルの組は\textbf{線型独立}\index{せんけいどくりつ@線型独立}であるという．

\begin{definition}{線型独立}{線型独立}
  $\bm{x}_1, \bm{x}_2, \dots, \bm{x}_n \in K^m$が\textbf{線型独立}であるとは，
  \[
    \lambda_1 \bm{x}_1 + \lambda_2 \bm{x}_2 + \cdots + \lambda_n \bm{x}_n = \bm{0}
  \]
  ならば，
  \[
    \lambda_1=\lambda_2=\cdots=\lambda_n = 0
  \]
  であることである．
  線型独立は\textbf{一次独立}\index{いちじどくりつ@一次独立}ともいう．
\end{definition}

\begin{prop}{線型独立の言い換え}{線型独立の言い換え}
  $\bm{x}_1,\bm{x}_2,\dots,\bm{x}_n\in K^m$が線型独立であることは，
  少なくとも一つが$0$でない任意の
  $\lambda_1,\lambda_2,\dots,\lambda_n\in K$に対して，
  \[
    \lambda_1\bm{x}_1+\lambda_2\bm{x}_2+\cdots
    +\lambda_n\bm{x}_n\ne\bm{0}
  \]
  となることと同値である．
\end{prop}

\begin{leftbar}
\begin{proof}
  \deref{線型独立}の条件の対偶を取れば明らかである．
\end{proof}
\end{leftbar}

線型独立でないベクトルの組は，\textbf{線型従属}\index{せんけいじゅうぞく@線型従属}であるという．

\begin{definition}{線型従属}{線型従属}
  $\bm{x}_1, \bm{x}_2, \dots, \bm{x}_n \in K^m$が線型独立でないとき，すなわち，少なくとも一つが$0$でない$\lambda_1, \lambda_2, \dots, \lambda_n \in K$が存在して
  \[
    \lambda_1 \bm{x}_1 + \lambda_2 \bm{x}_2 + \cdots + \lambda_n \bm{x}_n = \bm{0}
  \]
  が成り立つとき，$\bm{x}_1, \bm{x}_2, \dots, \bm{x}_n$は\textbf{線型従属}であるという．
  線型従属は\textbf{一次従属}\index{いちじじゅうぞく@一次従属}ともいう．
\end{definition}

\begin{example}{線型独立と線型従属}{線型独立と線型従属}
  $K^2$のベクトル
  \[
    \bm{e}_1 = \begin{pmatrix} 1 \\ 0 \end{pmatrix}, \quad
    \bm{e}_2 = \begin{pmatrix} 0 \\ 1 \end{pmatrix}
  \]
  は線型独立である．実際，
  \[
    \lambda_1 \bm{e}_1 + \lambda_2 \bm{e}_2 = \begin{pmatrix} \lambda_1 \\ \lambda_2 \end{pmatrix} = \bm{0}
  \]
  ならば$\lambda_1 = \lambda_2 = 0$となるからである．

  一方，$K^2$のベクトル
  \[
    \bm{x}_1 = \begin{pmatrix} 1 \\ 2 \end{pmatrix}, \quad
    \bm{x}_2 = \begin{pmatrix} 2 \\ 4 \end{pmatrix}
  \]
  は線型従属である．実際，$(\lambda_1, \lambda_2) = (2, -1) \ne (0, 0)$に対して
  \[
    2\bm{x}_1 + (-1)\bm{x}_2 = \bm{0}
  \]
  という非自明な線型関係が成り立つからである．
\end{example}

最後に，線型従属性の重要な言い換えを示そう．
これは「線型従属なベクトルの組では，どれかひとつのベクトルが残りのベクトルの線型結合として表せる」という主張であり，のちの章でも繰り返し用いられる．

\begin{prop}{線型従属の言い換え}{線型従属の言い換え}
  $n \ge 2$とする．
  $\bm{x}_1, \bm{x}_2, \dots, \bm{x}_n \in K^m$が線型従属であることは，ある$i$が存在して，$\bm{x}_i$が残りの$n - 1$個のベクトルの線型結合として表せることと同値である．
\end{prop}

\begin{leftbar}
\begin{proof}
  まず，$\bm{x}_1, \bm{x}_2, \dots, \bm{x}_n$が線型従属であるとする．
  このとき，少なくとも一つが$0$でないスカラー$\lambda_1, \lambda_2, \dots, \lambda_n \in K$が存在して
  \[
    \lambda_1 \bm{x}_1 + \lambda_2 \bm{x}_2 + \cdots + \lambda_n \bm{x}_n = \bm{0}
  \]
  が成り立つ．$\lambda_i \ne 0$となる$i$をとると，
  \[
    \bm{x}_i = \sum_{\substack{1 \le j \le n \\ j \ne i}} \left(-\frac{\lambda_j}{\lambda_i}\right) \bm{x}_j
  \]
  となり，$\bm{x}_i$は残りのベクトルの線型結合として表せる．

  逆に，ある$i$に対して
  \[
    \bm{x}_i = \sum_{\substack{1 \le j \le n \\ j \ne i}} \lambda_j \bm{x}_j
    \quad (\lambda_j \in K)
  \]
  と表せるとする．このとき，
  \[
    \lambda_1 \bm{x}_1 + \cdots + \lambda_{i-1} \bm{x}_{i-1} + (-1)\bm{x}_i + \lambda_{i+1} \bm{x}_{i+1} + \cdots + \lambda_n \bm{x}_n = \bm{0}
  \]
  が成り立つ．$\bm{x}_i$の係数は$-1 \ne 0$であるから，これは非自明な線型関係である．
  よって$\bm{x}_1, \bm{x}_2, \dots, \bm{x}_n$は線型従属である．

  以上の考察により，命題が示された．
\end{proof}
\end{leftbar}

すなわち，線型従属なベクトルの組には，残りのベクトルの線型結合として「すでに表せてしまう」ベクトルが存在するのである．

なお，与えられたベクトルの線型結合として表されるベクトル全体がなす集合は，のちの章において\textbf{線型包}として詳しく扱う．
そこでは，本章で定義した線型結合の概念が，数ベクトルに限らず一般の線型空間においてもそのまま通用することを見る．

\phantomsection
\subsection{線型独立性に関する基本補題}

線型独立性について，以後くり返し用いる二つの補題を用意しておこう．
いずれも線型結合と線型独立の定義のみから導かれる．

なお，これらの補題は，第3章において一般の線型空間に対してあらためて定式化され，証明される（\lemref{一般の取替補題}）．
ここでは，第2章までの議論で必要となる数ベクトル空間$K^m$の場合を先に示しておく．

\begin{lemma}{線型独立な組の拡大}{線型独立な組の拡大}
  $\bm{v}_1, \bm{v}_2, \dots, \bm{v}_k \in K^m$を線型独立とし，$\bm{w} \in K^m$は
  $\bm{v}_1, \bm{v}_2, \dots, \bm{v}_k$の線型結合として表せないとする．
  このとき，$\bm{v}_1, \bm{v}_2, \dots, \bm{v}_k, \bm{w}$は線型独立である．
\end{lemma}

\begin{leftbar}
\begin{proof}
  \[
    \lambda_1 \bm{v}_1 + \cdots + \lambda_k \bm{v}_k + \mu \bm{w} = \bm{0}
  \]
  とする．もし$\mu \ne 0$ならば，
  \[
    \bm{w} = \sum_{i=1}^{k} \left(-\frac{\lambda_i}{\mu}\right) \bm{v}_i
  \]
  となり，$\bm{w}$が線型結合として表せることになって仮定に反する．
  よって$\mu = 0$であり，このとき$\sum_{i=1}^{k} \lambda_i \bm{v}_i = \bm{0}$となるから，
  線型独立性より$\lambda_1 = \cdots = \lambda_k = 0$である．
  これが証明すべきことであった．
\end{proof}
\end{leftbar}

次の補題は，「線型独立なベクトルは，それを生成するベクトルの個数を超えて取れない」ことを述べている．

\begin{lemma}{取替補題}{取替補題}
  $\bm{v}_1, \bm{v}_2, \dots, \bm{v}_n \in K^m$とする．
  $\bm{w}_1, \bm{w}_2, \dots, \bm{w}_l \in K^m$がいずれも$\bm{v}_1, \bm{v}_2, \dots, \bm{v}_n$の線型結合として表され，
  かつ線型独立であるならば，
  \[
    l \le n
  \]
  が成り立つ．
\end{lemma}

\begin{leftbar}
\begin{proof}
  $k = 0, 1, \dots, l$に対して，次の主張$(\ast_k)$を$k$についての帰納法で示す．

  \begin{quote}
    $(\ast_k)$：$k \le n$であり，かつ$\bm{v}_1, \dots, \bm{v}_n$の番号を適当に付け替えることで，
    $\bm{v}_1, \dots, \bm{v}_n$の線型結合として表せるベクトルはすべて
    \[
      \bm{w}_1, \dots, \bm{w}_k, \bm{v}_{k+1}, \dots, \bm{v}_n
    \]
    の線型結合としても表せる，とできる．
  \end{quote}

  $k = 0$のときは主張の内容が空であり，成り立つ．

  $k < l$として$(\ast_k)$を仮定する．
  $\bm{w}_{k+1}$は$\bm{v}_1, \dots, \bm{v}_n$の線型結合であるから，帰納法の仮定より
  \begin{equation}
    \bm{w}_{k+1}
    = \sum_{i=1}^{k} a_i \bm{w}_i + \sum_{j=k+1}^{n} b_j \bm{v}_j
    \label{eq:取替補題の展開}
  \end{equation}
  と表せる．

  ここで，$b_{k+1} = \cdots = b_n = 0$であると仮定してみる
  （$k = n$のときは第$2$の和が空なので，自動的にこの場合にあたる）．
  このとき\eqref{eq:取替補題の展開}は
  \[
    a_1 \bm{w}_1 + \cdots + a_k \bm{w}_k + (-1)\bm{w}_{k+1} = \bm{0}
  \]
  と書き直せる．$\bm{w}_{k+1}$の係数は$-1 \ne 0$であるから，これは非自明な線型関係であり，
  $\bm{w}_1, \dots, \bm{w}_l$の線型独立性に矛盾する．

  よって$k < n$，すなわち$k + 1 \le n$であり，さらに$b_j \ne 0$となる$j$が存在する．
  残っている$\bm{v}_{k+1}, \dots, \bm{v}_n$の番号を付け替えて$b_{k+1} \ne 0$としてよい．
  このとき\eqref{eq:取替補題の展開}を$\bm{v}_{k+1}$について解けば，
  \[
    \bm{v}_{k+1}
    = \frac{1}{b_{k+1}}\left( \bm{w}_{k+1} - \sum_{i=1}^{k} a_i \bm{w}_i - \sum_{j=k+2}^{n} b_j \bm{v}_j \right)
  \]
  となり，$\bm{v}_{k+1}$は$\bm{w}_1, \dots, \bm{w}_{k+1}, \bm{v}_{k+2}, \dots, \bm{v}_n$の線型結合として表せる．
  したがって，$\bm{w}_1, \dots, \bm{w}_k, \bm{v}_{k+1}, \dots, \bm{v}_n$の線型結合として表せるベクトルは，
  すべて$\bm{w}_1, \dots, \bm{w}_{k+1}, \bm{v}_{k+2}, \dots, \bm{v}_n$の線型結合としても表せる．
  よって$(\ast_{k+1})$が成り立つ．

  以上より$(\ast_l)$が成り立ち，とくに$l \le n$である．
  これが証明すべきことであった．
\end{proof}
\end{leftbar}

\begin{cor}{$K^m$における線型独立なベクトルの個数}{Km における線型独立なベクトルの個数}
  $K^m$の線型独立なベクトルは，高々$m$個である．
  すなわち，$m + 1$個以上のベクトルは必ず線型従属である．
\end{cor}

\begin{leftbar}
\begin{proof}
  任意の$\bm{x} = (x_i) \in K^m$は
  \[
    \bm{x} = \sum_{i=1}^{m} x_i \bm{e}_i
  \]
  と標準基底ベクトルの線型結合で表せる．
  よって\lemref{取替補題}を$\bm{v}_i = \bm{e}_i$として適用すればよい．
  これが証明すべきことであった．
\end{proof}
\end{leftbar}

\begin{cor}{$K^n$の$n$個の線型独立なベクトル}{Kn の n 個の線型独立なベクトル}
  $\bm{a}_1, \bm{a}_2, \dots, \bm{a}_n \in K^n$が線型独立ならば，
  任意の$\bm{b} \in K^n$は$\bm{a}_1, \bm{a}_2, \dots, \bm{a}_n$の線型結合として表せる．
\end{cor}

\begin{leftbar}
\begin{proof}
  ある$\bm{b} \in K^n$が線型結合として表せないと仮定すると，\lemref{線型独立な組の拡大}より
  $\bm{a}_1, \dots, \bm{a}_n, \bm{b}$は線型独立である．
  これは$K^n$の$n + 1$個の線型独立なベクトルであり，\coref{Km における線型独立なベクトルの個数}に矛盾する．
  これが証明すべきことであった．
\end{proof}
\end{leftbar}

\phantomsection
\subsection{階数}

行列に対して，その列ベクトルがどれだけ「独立な情報」をもっているかを表す量を定める．

\begin{definition}{階数}{階数}
  $A$を$m \times n$行列とし，その列ベクトルを$\bm{a}_1, \bm{a}_2, \dots, \bm{a}_n \in K^m$とする．
  $\bm{a}_1, \bm{a}_2, \dots, \bm{a}_n$のうちから選んで線型独立となるベクトルの最大個数を，
  $A$の\textbf{階数}\index{かいすう@階数}\index{rank@$\rank$}といい，$\rank A$と表す．

  ただし，$0$個のベクトルの組は線型独立であると約束する．
  したがって$A = O$のとき$\rank A = 0$である．
\end{definition}

\begin{mycolumn}
  「線型独立な列の最大個数」という定義は，一見すると回りくどく感じられるかもしれない．
  実は，第3章で「次元」を定義したあとであれば，階数は「$A$の列ベクトルの線型結合として表せるベクトル全体（列空間とよばれる）の次元」と言い直すことができる．
  この列空間は，像$\Im T_A$にほかならない（\prref{線型写像と線型結合}のあとのコラムを参照）．
  すなわち階数は，第3章のことばでは「線型写像$T_A$の像の次元」$\dim \Im T_A$である（\coref{階数と像の次元}）．
  しかし本章の段階では次元がまだ定義されていないため，「次元」という語を使わずに同じ数を取り出す方法として，この定義を採用しているのである．
\end{mycolumn}

\begin{example}{階数}{階数の例}
  \[
    A = \begin{pmatrix} 1 & 2 & 0 \\ 2 & 4 & 0 \end{pmatrix}
  \]
  の列ベクトルは
  \[
    \bm{a}_1 = \begin{pmatrix} 1 \\ 2 \end{pmatrix}, \quad
    \bm{a}_2 = \begin{pmatrix} 2 \\ 4 \end{pmatrix}, \quad
    \bm{a}_3 = \begin{pmatrix} 0 \\ 0 \end{pmatrix}
  \]
  である．$\bm{a}_2 = 2\bm{a}_1$，$\bm{a}_3 = \bm{0}$であるから，
  どの$2$本を選んでも線型従属であり，一方$\bm{a}_1$の$1$本は線型独立である．
  よって$\rank A = 1$である．
\end{example}

階数の定義は列ベクトルによるものであったが，実は行ベクトルで定義しても同じ値が得られる．

\begin{theorem}{行階数と列階数の一致}{行階数と列階数の一致}
  $A$を$m \times n$行列とする．
  $A$の行ベクトルのうち線型独立なものの最大個数は，$\rank A$に等しい．
  したがって，
  \[
    \rank A^\mathsf{T} = \rank A
  \]
  が成り立つ．
\end{theorem}

\begin{leftbar}
\begin{proof}
  $\rank A = r$とし，$A$の行ベクトルのうち線型独立なものの最大個数を$r'$とする．

  まず$r' \le r$を示す．
  $A$の列ベクトル$\bm{a}_1, \dots, \bm{a}_n$のうち，線型独立な$r$本を選び，
  それらを$\bm{c}_1, \dots, \bm{c}_r$とする．
  各$j$について，$\bm{c}_1, \dots, \bm{c}_r, \bm{a}_j$は$r + 1$本であり，$r$の最大性から線型従属である．
  よって\lemref{線型独立な組の拡大}の対偶より，$\bm{a}_j$は$\bm{c}_1, \dots, \bm{c}_r$の線型結合として表せる：
  \[
    \bm{a}_j = \sum_{i=1}^{r} c_{ij} \bm{c}_i \quad (1 \le j \le n).
  \]

  そこで，$\bm{c}_1, \dots, \bm{c}_r$を列として並べた$m \times r$行列を$C$，
  $r \times n$行列$(c_{ij})$を$R$とおくと，上の式は$CR$の第$j$列が$\bm{a}_j$であることを意味するので，
  \[
    A = CR
  \]
  が成り立つ．
  行列の積の定義（\deref{行列の積}）より，$A$の第$k$行は
  \[
    (A \text{の第}k\text{行}) = \sum_{i=1}^{r} (C \text{の}(k,i)\text{成分}) \cdot (R \text{の第}i\text{行})
  \]
  とかける．すなわち，$A$の各行は$R$の$r$本の行の線型結合である．
  したがって，$A$の線型独立な行を$r'$本とれば，\lemref{取替補題}より$r' \le r$である．

  次に$r \le r'$を示す．
  前半の議論は任意の行列に対して通用する．すなわち，任意の行列$M$に対して
  \[
    (M\text{の線型独立な行の最大個数})
    \le
    (M\text{の線型独立な列の最大個数})
  \]
  が成り立つ．これを$M = A^\mathsf{T}$に適用する．
  $A^\mathsf{T}$の行は$A$の列であるから，左辺は$A$の線型独立な列の最大個数，すなわち$r$に等しい．
  また，$A^\mathsf{T}$の列は$A$の行であるから，右辺は$A$の線型独立な行の最大個数，すなわち$r'$に等しい．
  したがって$r \le r'$である．

  以上より$r = r'$であり，$A$の線型独立な行の最大個数は$\rank A$に等しい．
  さらに，$A^\mathsf{T}$の列は$A$の行であるから，
  \[
    \rank A^\mathsf{T} = r' = r = \rank A
  \]
  も成り立つ．
  これが証明すべきことであった．
\end{proof}
\end{leftbar}

\begin{cor}{階数の上限}{階数の上限}
  $m \times n$行列$A$に対して，
  \[
    \rank A \le \min\{m, n\}
  \]
  が成り立つ．
\end{cor}

\begin{leftbar}
\begin{proof}
  $A$の列は$n$本しかないので$\rank A \le n$である．
  また$A$の行は$m$本しかないので，\thref{行階数と列階数の一致}より$\rank A \le m$である．
  これが証明すべきことであった．
\end{proof}
\end{leftbar}

\begin{mycolumn}
  \thref{行階数と列階数の一致}は，一見すると当たり前ではない主張である．
  行の本数と列の本数は一般に異なるうえ，行の線型独立性と列の線型独立性は別々の条件だからである．
  それにもかかわらず両者が一致するという事実は，階数が行列の「かたち」ではなく本質的な情報を捉えていることを示している．

  なお，証明中に現れた分解$A = CR$は，「$m \times n$行列を，階数$r$の情報だけを使って$m \times r$行列と$r \times n$行列の積に分解する」ものであり，
  それ自体が応用上重要な意味をもつ．
\end{mycolumn}

\phantomsection
\subsection{階数の計算}

階数の定義はそのままでは計算に向かない．
実際には，行基本変形（\deref{行基本変形}）によって階数を求めることができる．

その根拠となるのが，次の観察である．

\begin{prop}{列の線型関係と斉次連立一次方程式}{列の線型関係と斉次連立方程式}
  $A$を$m \times n$行列とし，その列ベクトルを$\bm{a}_1, \dots, \bm{a}_n$とする．
  このとき，$\bm{x} = (x_j) \in K^n$に対して次は同値である：
  \begin{enumerate}[label=(\roman*),leftmargin=3.2em]
    \item $A\bm{x} = \bm{0}$である．
    \item $x_1 \bm{a}_1 + x_2 \bm{a}_2 + \cdots + x_n \bm{a}_n = \bm{0}$である．
  \end{enumerate}
  すなわち，$A$の列ベクトルの線型関係と，斉次連立一次方程式$A\bm{x} = \bm{0}$の解は同じものである．
\end{prop}

\begin{leftbar}
\begin{proof}
  \prref{線型写像と線型結合}のあとのコラムで見たとおり，
  \[
    A\bm{x} = x_1 \bm{a}_1 + x_2 \bm{a}_2 + \cdots + x_n \bm{a}_n
  \]
  が成り立つ．これが証明すべきことであった．
\end{proof}
\end{leftbar}

\begin{prop}{行基本変形と解集合}{行基本変形と解集合}
  $A$に行基本変形を施して得られる行列を$A'$とすると，
  \[
    \{\bm{x} \in K^n \mid A\bm{x} = \bm{0}\}
    = \{\bm{x} \in K^n \mid A'\bm{x} = \bm{0}\}
  \]
  が成り立つ．
\end{prop}

\begin{leftbar}
\begin{proof}
  $A\bm{x} = \bm{0}$の第$i$式を$f_i(\bm{x}) = 0$とかく．
  行基本変形の各操作で得られる新しい式は，いずれも$f_1, \dots, f_m$の線型結合である．
  よって$A\bm{x} = \bm{0}$を満たす$\bm{x}$は$A'\bm{x} = \bm{0}$も満たす．

  また，\deref{行基本変形}の$3$種類の操作はいずれも逆操作をもつ
  （第$i$行を$c$倍する操作の逆は$1/c$倍する操作，行の入れ替えの逆はもう一度入れ替える操作，
  第$i$行に第$j$行の$c$倍を加える操作の逆は$(-c)$倍を加える操作である）．
  同じ議論を逆操作に適用すれば，逆向きの包含も従う．
  これが証明すべきことであった．
\end{proof}
\end{leftbar}

\begin{theorem}{行基本変形と階数}{行基本変形と階数}
  行基本変形によって階数は変化しない．
\end{theorem}

\begin{leftbar}
\begin{proof}
  $A$に行基本変形を施して$A'$が得られたとする．
  列の番号$j_1 < j_2 < \cdots < j_k$を任意にとり，$A$の第$j_1, \dots, j_k$列だけを取り出した行列を$B$，
  $A'$の同じ列を取り出した行列を$B'$とする．
  $A$から$A'$への行基本変形は，$B$から$B'$への行基本変形をそのまま引き起こすから，
  \prref{行基本変形と解集合}より$B\bm{y} = \bm{0}$と$B'\bm{y} = \bm{0}$の解集合は一致する．

  \prref{列の線型関係と斉次連立方程式}より，
  $A$の第$j_1, \dots, j_k$列が線型独立であることは$B\bm{y} = \bm{0}$が自明な解しかもたないことと同値であり，
  $A'$についても同様である．
  よって，どの列の組についても，$A$で線型独立であることと$A'$で線型独立であることは一致する．
  線型独立な列の最大個数も一致するので，$\rank A = \rank A'$である．
  これが証明すべきことであった．
\end{proof}
\end{leftbar}

\begin{mycolumn}
  \thref{行基本変形と階数}により，階数は次のようにして求められる．
  行基本変形をくり返して$A$をできるだけ簡単な形に変形すると，
  最終的に，$0$でない各行の先頭の非零成分が右下に向かって階段状に並び，
  $0$のみからなる行が下側に集まった形にできる（さらに必要なら，各先頭成分を$1$にできる）．
  この形の行列では，$0$でない行はすべて線型独立であり，残りは$0$のみからなる行であるから，
  線型独立な行の最大個数は$0$でない行の本数に等しい．
  \thref{行階数と列階数の一致}より，これがそのまま階数である．

  \exref{階数の例}の行列$A$であれば，第$2$行に第$1$行の$(-2)$倍を加えて
  \[
    \begin{pmatrix} 1 & 2 & 0 \\ 2 & 4 & 0 \end{pmatrix}
    \longrightarrow
    \begin{pmatrix} 1 & 2 & 0 \\ 0 & 0 & 0 \end{pmatrix}
  \]
  となり，$0$でない行は$1$本であるから$\rank A = 1$と読み取れる．
\end{mycolumn}

\begin{mycolumn}
  階数は，のちに線型空間の次元を定義したあとで，
  「$A$の列ベクトルが張る線型部分空間（列空間）の次元」として捉え直される．
  \deref{階数}の「線型独立な列の最大個数」という定義は，その次元の値をあらかじめ直接に取り出したものにあたる．
\end{mycolumn}

ここまでで，線型代数の基礎となる概念についての議論を終えた．
次章ではこれらの概念を踏まえて，行列式を定義しよう．